%% =====================================================================
%%  Shadow Theory --- Paper 4
%%  Geometric Realization of Completed Source Relations
%%  Standalone LaTeX manuscript.  Compile with pdflatex (2x for refs).
%% =====================================================================
\documentclass[11pt,a4paper]{article}

\usepackage[T1]{fontenc}
\usepackage[utf8]{inputenc}
\usepackage{amsmath,amssymb,amsthm,mathrsfs}
\usepackage{mathtools}
\usepackage{bm}
\usepackage[margin=1in]{geometry}
\usepackage[expansion=false]{microtype}
\usepackage[colorlinks=true,linkcolor=blue,citecolor=blue,urlcolor=blue]{hyperref}
\usepackage{enumitem}

% ----- theorem environments -------------------------------------------
\theoremstyle{plain}
\newtheorem{theorem}{Theorem}[section]
\newtheorem{proposition}[theorem]{Proposition}
\newtheorem{lemma}[theorem]{Lemma}
\newtheorem{corollary}[theorem]{Corollary}

\theoremstyle{definition}
\newtheorem{definition}[theorem]{Definition}
\newtheorem{assumption}[theorem]{Assumption}
\newtheorem{example}[theorem]{Example}
\newtheorem{convention}[theorem]{Convention}

\theoremstyle{remark}
\newtheorem{remark}[theorem]{Remark}

% ----- macros ---------------------------------------------------------
\newcommand{\Real}{\mathbb{R}}
\newcommand{\Int}{\mathbb{Z}}
\newcommand{\Cx}{\mathbb{C}}
\newcommand{\gfrak}{\mathfrak{g}}
\newcommand{\Ggeo}{\mathfrak{G}_{\mathrm{geo}}}
\newcommand{\Cgeo}{\mathcal{C}_{\mathrm{geo}}}
\newcommand{\Sgeo}{S_{\mathrm{geo}}}
\newcommand{\Aut}{\operatorname{Aut}}
\newcommand{\AutB}{\operatorname{Aut}_{\mathcal{B}}}
\newcommand{\Gau}{\operatorname{Gau}}
\newcommand{\Diff}{\operatorname{Diff}}
\newcommand{\Hol}{\operatorname{Hol}}
\newcommand{\Ad}{\operatorname{Ad}}
\newcommand{\Tor}{\operatorname{Tor}}
\newcommand{\Imm}{\operatorname{Im}}
\newcommand{\id}{\mathrm{id}}
\newcommand{\dvol}{\mathrm{d}\!\operatorname{vol}_g}
\newcommand{\tr}{\operatorname{tr}}
\newcommand{\Egeo}{E_A^{\mathrm{geo}}}
\newcommand{\pgeo}{p_{\mathrm{geo}}}
\newcommand{\jgeo}{j_A^{\mathrm{geo}}}
\newcommand{\egeo}{e_A^{\mathrm{geo}}}
\newcommand{\Mflat}{\mathcal{M}_{\mathrm{flat}}}
\newcommand{\Lie}{\mathcal{L}}

\allowdisplaybreaks

% ======================================================================
\begin{document}

\title{\textbf{Geometric Realization of Completed Source Relations}\\[0.35em]
\large Descent, Orbit Spaces, Invariant Relations, and\\
Variational Response in Shadow Theory}
\author{Jeremy Rodgers\thanks{Website: \href{https://everythingequation.com}{everythingequation.com}}\\
\normalsize Independent Researcher}
\date{June 2026}
\maketitle

\begin{abstract}
The minimal completion constructed in Paper~3 is an abstract quotient determined by a
declared readout together with a set-indexed family of invariant source relations. This
paper identifies the conditions under which such a completion possesses a genuine
realization in differential geometry and classical field theory. The realization domain is
restricted to \emph{compatible} local metric, bundle, connection, matter, and boundary data:
\v{C}ech cocycle and overlap equations are used to construct global fields, and uniqueness is
asserted only up to bundle isomorphism, never up to a choice of coordinate chart or gauge
representative. The physical geometric source is then defined as the orbit space of admissible
global configurations under the groupoid of allowed principal-bundle isomorphisms covering
allowed diffeomorphisms.

Gauge- and diffeomorphism-invariant global relations---holonomy conjugacy classes, Wilson
observables, characteristic numbers, boundary and asymptotic data, and spectra of geometrically
defined operators with specified domains---descend to functions on this orbit space. Their joint
image with a declared Tier-1 readout realizes, on the restricted domain, the canonical minimal
completion of Paper~3. We prove a restricted geometric realization theorem, together with an
orbit reconstruction criterion characterizing when the completion map is injective. For flat
$U(1)$ connections on an oriented circle we prove that holonomy classifies gauge orbits
completely, yielding the exact realization $E_{\mathrm{hol}}\cong U(1)$, and we exhibit a
strictly coarser spectral completion $E_{\mathrm{spec}}\cong[-1,1]$ arising from a charged
scalar relation.

The final sections embed completed relation data into variational dynamics. A gauge- and
diffeomorphism-invariant relation-dependent action yields correctly normalized stress-energy,
gauge-current, and matter responses; the Einstein, Yang--Mills, and matter equations follow by
functional variation, and the Ward--Noether identities associated with diffeomorphism and gauge
invariance enforce covariant conservation. Completed source relations therefore enter Tier-1
field equations through invariant geometric structure and derived response, not through a
choice of coordinates or gauge. No claim is made that every abstract completion is
geometrically realizable; the realization boundary is precisely the content of the results
below.
\end{abstract}

\tableofcontents

% ======================================================================
\section{Introduction and scientific position}
\label{sec:intro}

\subsection{Position in the sequence}

Paper~1 \cite{ST1} established source--readout non-equivalence and identified the descent
obstructions that prevent a readout from being inverted. Paper~2 \cite{ST2} characterized when
the information lost by a readout obstructs a nominated target, and what any target-sufficient
extension is obliged to distinguish. Paper~3 \cite{ST3} constructed, for a set-indexed family
$\{R_\alpha\}_{\alpha\in A}$ of invariant source relations and a declared readout
$p\colon S\to T$, the canonical minimal completion
\begin{equation}
\label{eq:paper3-completion}
E_A \;=\; \Imm\bigl(p,R_A\bigr)
\;\subseteq\; T\times\prod_{\alpha\in A}V_\alpha ,
\qquad
R_A=(R_\alpha)_{\alpha\in A},
\end{equation}
together with its quotient realization $S\twoheadrightarrow E_A$ and its terminal universal
property among surjective extensions through which $p$ and every $R_\alpha$ factor.

Construction \eqref{eq:paper3-completion} is purely set-theoretic. It is indifferent to whether
the relation values $(t,(v_\alpha))$ can be produced by any geometry at all. Paper~4 supplies
exactly the missing ingredient:

\begin{quote}
\emph{When do the abstract relation coordinates of Paper~3 arise from compatible global
metrics, bundles, connections, matter fields, boundary data, and gauge-invariant observables?}
\end{quote}

The answer has the structure of a \emph{realization boundary}. Not every point of the Cartesian
product $T\times\prod_\alpha V_\alpha$ is attained; the attainable points are precisely the
image of a geometric evaluation map, and that image is constrained by descent (\v{C}ech
cocycle and overlap conditions), by invariance (the relations must be constant on orbits of the
allowed symmetry groupoid), and by the admissibility class chosen for the model (regularity,
signature, causal structure, boundary and asymptotic conditions). The abstract completion is
then realized \emph{relative to} the invariant geometric relations declared, and only on that
restricted domain.

\subsection{Progression of the paper}

The paper follows a single line of development. Compatible local geometric data glue to global
fields (Section~\ref{sec:descent}, Theorem~\ref{thm:descent}). Physical configurations are
reduced by allowed bundle isomorphisms and diffeomorphisms (Section~\ref{sec:source},
Definition~\ref{def:sgeo}). Invariant global geometric relations become well-defined functions
on the reduced orbit space (Section~\ref{sec:invariants}). Their joint image with a declared
readout realizes the canonical completion of Paper~3 (Section~\ref{sec:realization},
Theorem~\ref{thm:restricted-realization}). Flat $U(1)$ connections on an oriented circle provide
an exact classification example (Section~\ref{sec:u1}, Theorem~\ref{thm:hol-class}). Invariant
relation-dependent actions generate stress-energy, currents, and matter response through
functional variation (Section~\ref{sec:variational}, Theorem~\ref{thm:field-eqs}). Gauge and
diffeomorphism invariance impose the corresponding Ward--Noether identities
(Section~\ref{sec:noether}, Theorem~\ref{thm:cons} and Proposition~\ref{prop:gauge-current}).

\subsection{Scope}
\label{sec:scope}

The scope restrictions recorded here are load-bearing and are used explicitly in the statements
of the theorems below. It is not claimed that every abstract completion of Paper~3 is
geometrically realizable; realizability is a nontrivial condition, and
Corollary~\ref{cor:realizability-boundary} exhibits the obstruction. It is not claimed that any
finite list of curvature scalars separates arbitrary Lorentzian geometries; separation is a
hypothesis to be established on the restricted domain of a given model, and
Remark~\ref{rem:no-finite-separation} records why no general result of that kind is available.
Gauge potentials and coordinate representatives are never treated as physical states: all
realization relations are functions on the orbit space $\Sgeo$. A loop observable is not treated
as diffeomorphism invariant unless the loop itself is included in the configuration data and
transformed with it, or the observable is formed by an invariant construction over loops, or the
declared boundary data pins the relevant homology class; Section~\ref{sec:diff-loops} makes these
alternatives precise and requires a model to declare which it uses. Global relation coordinates
are not treated as independent dynamical fields unless a model explicitly promotes them to
fields. No new physical interaction or empirical claim is introduced, and no claim is made about
the physical status of our universe: the relation-dependent action of
Section~\ref{sec:variational} is a template for admissible couplings, and its consequences are
theorems about that template, not predictions. Finally, the projected-dynamics and memory
machinery of Paper~5, and the Randall--Sundrum construction of Paper~7, are not used and are not
assumed; where the present results terminate in an open dynamical question, this is stated as
such (Section~\ref{sec:tier1}).

% ======================================================================
\section{The geometric source domain}
\label{sec:source}

\subsection{Configurations}

\begin{convention}[Smoothness, signature, and index conventions]
\label{conv:basic}
Throughout the general theory, $M$ is a smooth, connected, paracompact, oriented four-manifold,
possibly with boundary $\partial M$. The choice $\dim M=4$ is made for definiteness in the physical
application and is used nowhere essentially: Sections~\ref{sec:source}--\ref{sec:realization} hold
verbatim in any dimension. Section~\ref{sec:u1} is a deliberate one-dimensional specialization,
declared as such in Convention~\ref{conv:circle-scope}, and the dimensional stipulation of the
present convention is suspended there. All fields are smooth unless a weaker regularity class is explicitly
declared as part of the admissibility data. A Lorentzian metric has signature $(-,+,+,+)$; a
\emph{time-oriented} Lorentzian metric is a pair $(g,\mathfrak{t})$ consisting of a Lorentzian
metric $g$ together with a choice $\mathfrak{t}$ of one of the two continuous global
orientations of the timelike cone bundle of $g$ (equivalently, the choice of a connected
component of the set of $g$-timelike vector fields, when one exists). Not every Lorentzian
metric admits such a choice; admitting one is a nontrivial global condition on $g$, and it is
imposed below as part of admissibility rather than derived
(see Remark~\ref{rem:time-orientation}). Greek indices run over $0,1,2,3$ and are raised and
lowered with $g$. $G$ denotes a finite-dimensional Lie group with Lie algebra $\gfrak$, equipped
with a fixed $\Ad$-invariant nondegenerate bilinear form $\langle\cdot,\cdot\rangle$ on $\gfrak$
(for compact semisimple $G$ one may take the negative Killing form; for $G=U(1)$ the standard
normalization on $\mathfrak{u}(1)\cong i\Real$). Finally $\rho\colon G\to \mathrm{GL}(W)$ is a
fixed finite-dimensional representation; $\rho$ is \emph{not} assumed unitary, and where
unitarity is needed it is hypothesized explicitly.
\end{convention}

\begin{definition}[Global geometric configuration]
\label{def:config}
A \emph{global geometric configuration} is a tuple
\begin{equation}
\label{eq:config}
\Psi=\bigl(M,\;g,\;P,\;A,\;\Phi,\;\mathcal{B}\bigr),
\end{equation}
in which
\begin{enumerate}[label=(\alph*),leftmargin=2.2em]
\item $g$ is a Lorentzian metric on $M$;
\item $\pi\colon P\to M$ is a principal $G$-bundle;
\item $A\in\Omega^1(P,\gfrak)$ is a principal connection on $P$
      \cite{KobayashiNomizu1,Bleecker};
\item $\Phi$ is a section of the associated bundle $P\times_\rho W\to M$;
\item $\mathcal{B}$ denotes the boundary, asymptotic, causal, and global data required to make
      the model well posed. Among these data we \emph{always} include a choice of time
      orientation $\mathfrak{t}$ for $g$; the remaining components of $\mathcal{B}$ are
      model-specific and may comprise, for instance, fixed induced data on $\partial M$, a decay
      class at infinity, a fixed conformal boundary structure, a global hyperbolicity
      requirement, a spin structure, or a fixed topological sector.
\end{enumerate}
\end{definition}

\begin{remark}[Where time orientation lives]
\label{rem:time-orientation}
Time orientation is a property of the pair $(M,g)$, not of the bare smooth manifold $M$: it is
meaningless to speak of $M$ as ``time-orientable'' independently of a Lorentzian metric, since
the timelike cone bundle whose components are being selected is defined by $g$. Accordingly the
choice $\mathfrak{t}$ is carried in $\mathcal{B}$, and the arrows of the redundancy groupoid are
required below to preserve it. This has a consequence for the descent theorem, recorded
explicitly in Theorem~\ref{thm:descent} and Remark~\ref{rem:descent-time}: descent reconstructs
the \emph{metric}, but membership of the reconstructed $(g,\mathfrak{t})$ in a time-oriented
admissibility class is an additional \emph{global} condition on the glued object, which either
must be imposed after the fact or must be built into the descent data as compatible local time
orientations.
\end{remark}

\begin{definition}[Admissibility class]
\label{def:admissible}
An \emph{admissibility class} $\Cgeo$ is a specified class of tuples \eqref{eq:config}
satisfying declared conditions on regularity, signature, causal structure, boundary and
asymptotic behaviour, and topological sector, and in which every $g$ is time-orientable with a
distinguished time orientation $\mathfrak{t}\in\mathcal{B}$. The class $\Cgeo$ is part of the
data of the model. All subsequent statements are relative to a fixed $\Cgeo$.
\end{definition}

The insistence that $\Cgeo$ be declared is not pedantry. Every separation, realization, and
conservation statement below is sensitive to it: a relation family may separate orbits in one
admissibility class and fail to separate them in a larger one, and a Ward identity may hold with
compactly supported variations and fail without a boundary term.

\subsection{The redundancy groupoid}

Gauge and coordinate freedom in a bundle-theoretic setting is not, in general, a group acting
on a fixed configuration space: different admissible configurations may carry non-isomorphic
bundles, and the diffeomorphisms that lift to a given bundle need not lift to another. The
correct structure is a groupoid.

\begin{definition}[Ambient allowed automorphisms]
\label{def:ambient-aut}
Fix a principal $G$-bundle $P\to M$. Let $\Diff_{\mathcal{B}}(M)$ denote the group of
diffeomorphisms of $M$ preserving the orientation of $M$ and the non-metric structural data in
$\mathcal{B}$ (boundary, asymptotic regime, topological sector), and let
\begin{equation}
\label{eq:AutB}
\AutB(P)
:=\bigl\{\hat F\in\Aut(P)\;:\;\hat F \text{ covers some } f\in\Diff_{\mathcal{B}}(M)\bigr\}
\end{equation}
be the group of $G$-equivariant diffeomorphisms of $P$ covering such diffeomorphisms. Define
\begin{equation}
\label{eq:DiffPB}
\Diff_{P,\mathcal{B}}(M)\;:=\;\text{image of }\ \AutB(P)\ \text{ in }\ \Diff_{\mathcal{B}}(M),
\end{equation}
that is, the subgroup of allowed base diffeomorphisms that admit an equivariant lift to $P$.
\end{definition}

\begin{proposition}[Structure of the ambient group]
\label{prop:aut-sequence}
With the definitions \eqref{eq:AutB}--\eqref{eq:DiffPB} there is a short exact sequence of groups
\begin{equation}
\label{eq:aut-sequence}
1\;\longrightarrow\;\Gau(P)\;\longrightarrow\;\AutB(P)\;\longrightarrow\;
\Diff_{P,\mathcal{B}}(M)\;\longrightarrow\;1 ,
\end{equation}
where
\[
\Gau(P)\;=\;\bigl\{\hat F\in\Aut(P):\hat F\ \text{covers}\ \id_M\bigr\}
\;\cong\;\Gamma\bigl(P\times_{\mathrm{conj}}G\bigr)
\]
is the gauge group. The sequence is exact by construction: $\Diff_{P,\mathcal{B}}(M)$ is
\emph{defined} as the image of the projection, and the kernel of that projection is exactly the
set of automorphisms covering the identity. The sequence need not split.
\end{proposition}

\begin{proof}
Exactness at $\Gau(P)$ and at $\Diff_{P,\mathcal{B}}(M)$ is immediate from the definitions of
kernel and image. That the extension need not split is a statement about bundle topology: a
splitting would furnish, for each allowed $f\in\Diff_{P,\mathcal{B}}(M)$, an equivariant lift
$\sigma(f)\in\AutB(P)$ depending on $f$ homomorphically, and there is no reason for such a
coherent choice to exist. Consequently $\AutB(P)$ is not in general a semidirect product
$\Gau(P)\rtimes\Diff_{P,\mathcal{B}}(M)$, and nothing below assumes it is.
\end{proof}

\begin{remark}[The ambient group is not the isotropy group]
\label{rem:isotropy}
Proposition~\ref{prop:aut-sequence} describes the \emph{ambient} group of allowed bundle
automorphisms over a fixed $P$. It does \emph{not} describe the isotropy group of the groupoid at
a configuration. The isotropy (stabilizer) group at
$\Psi=(M,g,P,A,\Phi,\mathcal{B})$ is the subgroup
\begin{equation}
\label{eq:stabilizer}
\begin{aligned}
\Aut(\Psi)
:=\bigl\{\hat F\in\AutB(P)\ \text{covering some}\ f\ :\ \ 
&f_*g=g,\quad f_*\mathfrak{t}=\mathfrak{t},\quad (\hat F^{-1})^*A=A,\\[-0.15em]
&\hat F_\rho\circ\Phi=\Phi\circ f,\quad f_*\mathcal{B}=\mathcal{B}
\bigr\}\;\le\;\AutB(P),
\end{aligned}
\end{equation}
that is, the stabilizer of the \emph{whole tuple}: it preserves the metric, the chosen time
orientation and orientation, the connection, the matter section, and the boundary data.
$\Aut(\Psi)$ is typically far smaller than $\AutB(P)$: for a generic $(g,A,\Phi)$ it is a small or
even trivial group, its base part being contained in the isometry group of $g$. In particular the
exact sequence \eqref{eq:aut-sequence} is a statement about $\AutB(P)$ and not about
$\Aut(\Psi)$.
\end{remark}

\begin{definition}[Redundancy groupoid]
\label{def:groupoid}
Let $\Ggeo$ be the groupoid whose objects are the elements of $\Cgeo$. Given two objects
\[
\Psi=(M,g,P,A,\Phi,\mathcal{B}),
\qquad
\Psi'=(M,g',P',A',\Phi',\mathcal{B}'),
\]
the arrows $F\colon\Psi\to\Psi'$ are the pairs $F=(\hat F,f)$ consisting of
\begin{itemize}[leftmargin=1.6em]
\item an allowed diffeomorphism $f\in\Diff_{\mathcal{B}}(M)$, and
\item a $G$-equivariant bundle isomorphism $\hat F\colon P\to P'$ covering $f$,
\end{itemize}
subject to
\begin{equation}
\label{eq:arrow-conditions}
f_*g=g',
\qquad
f_*\mathfrak{t}=\mathfrak{t}',
\qquad
(\hat F^{-1})^{*}A = A',
\qquad
\hat F_{\rho}\circ\Phi = \Phi'\circ f,
\qquad
f_*\mathcal{B}=\mathcal{B}',
\end{equation}
where $\hat F_\rho$ is the induced isomorphism $P\times_\rho W\to P'\times_\rho W$. The second
condition is the requirement that arrows preserve the chosen time orientation; it is not implied
by the first, since $f_*g=g'$ alone permits $f$ to exchange the two components of the timelike
cone bundle. Composition is composition of bundle maps; identities are $(\id_P,\id_M)$. The
isotropy group of $\Ggeo$ at $\Psi$ is $\Aut(\Psi)$ of \eqref{eq:stabilizer}.
\end{definition}

The groupoid formulation---rather than a single group acting on a single configuration
space---is forced by two features. First, the arrows relate objects carrying \emph{different}
bundle representatives, which is exactly what is needed when descent produces a bundle only up to
isomorphism (Theorem~\ref{thm:descent}). Second, isotropy jumps from configuration to
configuration, since $\Aut(\Psi)$ depends on $\Psi$; there is no ambient group acting with
constant stabilizer.

\begin{definition}[Physical equivalence and the reduced source]
\label{def:sgeo}
Write
\begin{equation}
\label{eq:phys-equiv}
\Psi\simeq_{\mathrm{geo}}\Psi'
\quad\Longleftrightarrow\quad
\exists\,F\colon\Psi\to\Psi' \text{ in } \Ggeo .
\end{equation}
This is an equivalence relation on $\Cgeo$ because $\Ggeo$ is a groupoid. The \emph{reduced
geometric source} is the orbit space
\begin{equation}
\label{eq:sgeo}
\Sgeo \;:=\; \Cgeo/\!\simeq_{\mathrm{geo}} ,
\end{equation}
with quotient map $q\colon\Cgeo\to\Sgeo$, $q(\Psi)=[\Psi]$.
\end{definition}

\begin{remark}[Set-theoretic hygiene]
\label{rem:hygiene}
$\Sgeo$ is used below only as a set (equivalently, as the set of isomorphism classes of objects
of the essentially small groupoid $\Ggeo$ obtained after fixing a set of bundle
representatives). No smooth or stratified structure on $\Sgeo$ is used or assumed; the orbit
space of a gauge groupoid is in general not a manifold---precisely because isotropy jumps, as
Remark~\ref{rem:isotropy} notes---and the arguments of
Sections~\ref{sec:realization}--\ref{sec:u1} do not require it to be. Where a topology is
needed, as for the moduli space $\Mflat(S^1,U(1))$ in Section~\ref{sec:u1}, it is supplied and
identified explicitly.
\end{remark}

\begin{convention}[Status of relation coordinates]
\label{conv:relation-status}
A relation coordinate $r_\alpha$ is a \emph{function on $\Sgeo$}, not an independent field on
$M$. It is not varied independently in Section~\ref{sec:variational} unless a model explicitly
promotes it to a field by adjoining it to the configuration tuple \eqref{eq:config} and
enlarging $\Cgeo$ accordingly. Where $r_\alpha$ is a functional of $(g,A,\Phi)$ its variation is
computed by the functional chain rule of Section~\ref{sec:chain-rule}; where it labels a fixed
topological or boundary sector it is held fixed under variations \emph{within} that sector. Both
cases are handled in Section~\ref{sec:variational}; they are not conflated.
\end{convention}

Every realization relation used in this paper is, by construction, a function on $\Sgeo$ and
never a function of an arbitrary coordinate chart or gauge potential.

% ======================================================================
\section{Local data and global descent}
\label{sec:descent}

We now establish the first leg of the progression: compatible local data glue to global fields.
This is the theorem that fixes the \emph{domain} of geometric realization. The material is
classical \cite{KobayashiNomizu1,Steenrod,Bleecker}; what is new is only the use to which it is
put and the care with which the invariance and admissibility hypotheses are tracked.

\subsection{Local data and compatibility equations}

Let $\{U_i\}_{i\in I}$ be an open cover of $M$. Suppose we are given, for each $i\in I$,
\begin{equation}
\label{eq:local-data}
g_i\in\Gamma\bigl(S^2T^*U_i\bigr),
\qquad
A_i\in\Omega^1(U_i,\gfrak),
\qquad
\Phi_i\in C^\infty(U_i,W),
\end{equation}
where each $g_i$ is an \emph{intrinsic} smooth symmetric covariant $2$-tensor field on the open
set $U_i\subseteq M$, of Lorentzian signature at every point; together with smooth transition
functions
\begin{equation}
\label{eq:transition}
h_{ij}\colon U_i\cap U_j\longrightarrow G .
\end{equation}
If in addition each $U_i$ is the domain of a chart $\psi_i\colon U_i\to\psi_i(U_i)\subseteq\Real^4$,
we may pass to \emph{coordinate components}. These live on the coordinate images, not on $M$:
setting
\begin{equation}
\label{eq:components}
\bigl(\check g_i\bigr)_{ab}
:=\bigl((\psi_i^{-1})^*g_i\bigr)_{ab}
\quad\text{on}\quad \psi_i(U_i)\subseteq\Real^4 ,
\end{equation}
and writing $\psi_{ij}:=\psi_i\circ\psi_j^{-1}$ for the coordinate transition map on
$\psi_j(U_i\cap U_j)$, the components of a \emph{single} intrinsic tensor in two charts are
related by the usual pullback law
\begin{equation}
\label{eq:component-law}
\check g_j \;=\; \psi_{ij}^{\,*}\,\check g_i
\qquad\text{on}\qquad \psi_j(U_i\cap U_j),
\qquad\text{i.e.}\qquad
\bigl(\check g_j\bigr)_{cd}
=\frac{\partial \psi_{ij}^{\,a}}{\partial x^{c}}\,
 \frac{\partial \psi_{ij}^{\,b}}{\partial x^{d}}\,
 \bigl(\check g_i\bigr)_{ab}\circ\psi_{ij}.
\end{equation}
Equation \eqref{eq:component-law} is a statement about \emph{components on coordinate images}. It
is distinct from the gluing condition \eqref{eq:metric-overlap} below, which is a relation among the
intrinsic tensors $g_i$ on their overlaps.

\begin{definition}[Compatible local geometric data]
\label{def:compatible}
The collection $\{g_i,h_{ij},A_i,\Phi_i\}$ is \emph{compatible} if, on every nonempty overlap
(respectively triple overlap):
\begin{align}
\label{eq:cocycle}
h_{ii}&=1,
&
h_{ji}&=h_{ij}^{-1},
&
h_{ij}h_{jk}&=h_{ik} ,
\\[0.3em]
\label{eq:metric-overlap}
g_i\big|_{U_i\cap U_j} &= g_j\big|_{U_i\cap U_j} ,
\\[0.3em]
\label{eq:connection-overlap}
A_j &= \Ad_{h_{ij}^{-1}}A_i + h_{ij}^{*}\theta_G ,
\\[0.3em]
\label{eq:matter-overlap}
\Phi_j &= \rho\bigl(h_{ij}^{-1}\bigr)\Phi_i .
\end{align}
Here \eqref{eq:cocycle} is the \v{C}ech cocycle condition; \eqref{eq:metric-overlap} is the
gluing condition (agreement of restrictions) for a section of the sheaf $S^2T^*M$; and \eqref{eq:connection-overlap}, \eqref{eq:matter-overlap} are
the compatibility conditions for local connection forms and local matter representatives,
$\theta_G$ denoting the left Maurer--Cartan form on $G$. For a matrix group
$G\subseteq\mathrm{GL}(n)$, equation \eqref{eq:connection-overlap} takes the familiar form
\[
A_j = h_{ij}^{-1}A_i\,h_{ij} + h_{ij}^{-1}\,\mathrm{d}h_{ij} .
\]
\end{definition}

\begin{remark}[\v{C}ech $H^1$ for nonabelian $G$]
\label{rem:cech}
A cocycle $\{h_{ij}\}$ satisfying \eqref{eq:cocycle} represents a class in
$\check H^1(\{U_i\},\underline{G})$, where $\underline{G}$ is the sheaf of smooth $G$-valued
functions. For nonabelian $G$ this object is a \emph{pointed set} of cocycle classes---the
distinguished point being the class of the trivial cocycle---and not a group; two cocycles are
equivalent when they are cohomologous in the sense
$h'_{ij}=\lambda_i^{-1}h_{ij}\lambda_j$ for smooth $\lambda_i\colon U_i\to G$. Passing to the
direct limit over refinements gives the pointed set $\check H^1(M,\underline{G})$, which
classifies principal $G$-bundles over $M$ up to isomorphism \cite{Steenrod}. Only for abelian $G$
does this pointed set carry a natural group structure; we use it below only as a pointed set.
\end{remark}

\begin{assumption}[Regularity of the descent data]
\label{ass:descent-reg}
The cover $\{U_i\}$ is locally finite; the $g_i$, $A_i$, $\Phi_i$, $h_{ij}$ are smooth on their
domains; and each $g_i$ has the signature fixed in Convention~\ref{conv:basic}. If a weaker
regularity class is used in $\Cgeo$, the same statements hold with ``smooth'' replaced uniformly
by that class, provided the class is closed under restriction and under the gauge transformations
appearing in \eqref{eq:connection-overlap}--\eqref{eq:matter-overlap}.
\end{assumption}

\subsection{The descent theorem}

\begin{theorem}[Geometric descent]
\label{thm:descent}
Let $\{g_i,h_{ij},A_i,\Phi_i\}$ be compatible local geometric data in the sense of
Definition~\ref{def:compatible}, satisfying Assumption~\ref{ass:descent-reg}. Then these data
determine:
\begin{enumerate}[label=(\arabic*),leftmargin=2.2em]
\item a global Lorentzian metric $g$ on $M$ with $g|_{U_i}=g_i$;
\item a principal $G$-bundle $\pi\colon P\to M$ with local trivializations over $\{U_i\}$ whose
      transition functions are $\{h_{ij}\}$;
\item a global principal connection $A$ on $P$ whose local representatives (pullbacks along the
      canonical local sections $s_i$) are the $A_i$;
\item a global section $\Phi$ of $P\times_\rho W$ whose local representatives are the $\Phi_i$.
\end{enumerate}
The resulting global tuple $(g,P,A,\Phi)$ is unique up to the bundle isomorphisms induced by
change of local trivialization. Conversely, every global tuple $(g,P,A,\Phi)$ over $M$ arises in
this way from compatible local data, for any cover trivializing $P$.

The theorem reconstructs the metric $g$; it does \emph{not} by itself produce a time orientation.
Membership of $(M,g)$ in the time-oriented admissibility class $\Cgeo$ of
Definition~\ref{def:admissible} is an additional \emph{global} condition on the glued metric,
discussed in Remark~\ref{rem:descent-time}.
\end{theorem}

\begin{proof}
\emph{(1) The metric.}
Condition \eqref{eq:metric-overlap} says exactly that the intrinsic tensors $g_i$ and $g_j$
restrict to the same tensor field on $U_i\cap U_j$. Since $\Gamma(S^2T^*-)$ is a sheaf on $M$, the
gluing axiom yields a unique global section $g\in\Gamma(S^2T^*M)$ with $g|_{U_i}=g_i$ for all $i$.
Nondegeneracy and signature are pointwise properties of the symmetric bilinear form $g_x$ on
$T_xM$; every $x\in M$ lies in some $U_i$, where $g_x=(g_i)_x$ is Lorentzian by hypothesis. Hence
$g$ is Lorentzian. (If one prefers to present the data in coordinate components, the equivalent
hypothesis is \eqref{eq:component-law}, which expresses that the $\check g_i$ are the component
arrays of a single intrinsic tensor; the invariance of signature under the linear isomorphism
$\mathrm{d}\psi_{ij}(x)$ is then Sylvester's law of inertia.)

\emph{(2) The bundle.}
Form the disjoint union $\coprod_{i\in I}U_i\times G\times\{i\}$ and impose
\begin{equation}
\label{eq:bundle-glue}
(x,a,i)\;\sim\;\bigl(x,\;h_{ji}(x)\,a,\;j\bigr),
\qquad x\in U_i\cap U_j,\; a\in G .
\end{equation}
Reflexivity follows from $h_{ii}=1$, symmetry from $h_{ij}=h_{ji}^{-1}$, and transitivity from
$h_{kj}h_{ji}=h_{ki}$; hence \eqref{eq:bundle-glue} is an equivalence relation. Let $P$ be the
quotient with the quotient topology, $\pi[(x,a,i)]:=x$.

We record the trivializations rigorously. Fix $i$ and $x\in U_i$. \emph{Every equivalence class
lying over $x$ has exactly one representative whose third entry is $i$}: existence, because if a
class contains $(x,b,j)$ then it contains $(x,h_{ij}(x)b,i)$; uniqueness, because
$(x,a,i)\sim(x,a',i)$ forces $a'=h_{ii}(x)a=a$. Hence the map
\begin{equation}
\label{eq:triv}
\varphi_i\colon \pi^{-1}(U_i)\longrightarrow U_i\times G,
\qquad
\varphi_i\bigl([(x,a,i)]\bigr)=(x,a)
\end{equation}
is a well-defined bijection, with inverse $\varphi_i^{-1}(x,a)=[(x,a,i)]$; it is a homeomorphism
by definition of the quotient topology, and we give $P$ the smooth structure for which the
$\varphi_i$ are diffeomorphisms (consistent, since the change of trivialization
$\varphi_i\circ\varphi_j^{-1}(x,a)=(x,h_{ij}(x)a)$ is smooth). The right $G$-action
$[(x,a,i)]\cdot b:=[(x,ab,i)]$ is well defined, because $h_{ji}$ acts by \emph{left}
multiplication and hence commutes with right multiplication; it is free and transitive on fibres,
and each $\varphi_i$ is $G$-equivariant. Thus $\pi\colon P\to M$ is a principal $G$-bundle with
transition functions $h_{ij}$.

\emph{(3) The connection.}
Let $s_i\colon U_i\to P$ be the canonical local sections $s_i(x):=\varphi_i^{-1}(x,1)$, so that
$s_j=s_i\cdot h_{ij}$ on $U_i\cap U_j$. A principal connection $A\in\Omega^1(P,\gfrak)$ is
determined by its local representatives $A_i:=s_i^*A$, and a family $\{A_i\}$ arises from a
(necessarily unique) global $A$ if and only if it obeys the law induced by $s_j=s_i\cdot h_{ij}$,
namely
\[
s_j^*A \;=\; \Ad_{h_{ij}^{-1}}\bigl(s_i^*A\bigr)+h_{ij}^*\theta_G,
\]
which is precisely \eqref{eq:connection-overlap}. Given \eqref{eq:connection-overlap}, define $A$
over $U_i$ by the pullback identity
\begin{equation}
\label{eq:conn-reconstruct}
\bigl(\varphi_i^{-1}\bigr)^{*}A
\;=\;
\Ad_{a^{-1}}\bigl(\pi_U^{*}A_i\bigr)\;+\;\theta_G
\qquad\text{on}\qquad U_i\times G,
\end{equation}
where $\pi_U\colon U_i\times G\to U_i$ and the projection $U_i\times G\to G$ are understood, $a$
denotes the $G$-coordinate, and $\theta_G$ is pulled back along that projection. Formula
\eqref{eq:conn-reconstruct} is an equality of $\gfrak$-valued $1$-forms on the manifold
$U_i\times G$; it is not an evaluation at a point. Two such local definitions agree on
$\pi^{-1}(U_i\cap U_j)$: transporting \eqref{eq:conn-reconstruct} for $j$ through
$\varphi_i\circ\varphi_j^{-1}(x,a)=(x,h_{ij}(x)a)$ and using the standard identity
$(L_{h}\circ)^*\theta_G=\theta_G$ together with the Maurer--Cartan expansion of
$\mathrm{d}(h_{ij}a)$ reproduces \eqref{eq:conn-reconstruct} for $i$ exactly when
\eqref{eq:connection-overlap} holds. The resulting $1$-form is $G$-equivariant,
$R_b^*A=\Ad_{b^{-1}}A$, and reproduces the fundamental vector fields, $A(\zeta_X)=X$; hence it is
a principal connection \cite{KobayashiNomizu1}.

\emph{(4) The matter section.}
Sections of $P\times_\rho W$ correspond bijectively to $\rho$-equivariant maps
$\tilde\Phi\colon P\to W$ with $\tilde\Phi(p\cdot a)=\rho(a^{-1})\tilde\Phi(p)$, and such maps
correspond to families $\Phi_i:=\tilde\Phi\circ s_i$ obeying
$\Phi_j=\tilde\Phi\circ(s_i\cdot h_{ij})=\rho(h_{ij}^{-1})\Phi_i$, which is
\eqref{eq:matter-overlap}. Hence the $\Phi_i$ glue to a unique global section $\Phi$.

\emph{Uniqueness up to isomorphism.}
Suppose two systems of local trivializations over a common cover (or over covers admitting a
common refinement, after passing to it) produce compatible data $\{g_i,h_{ij},A_i,\Phi_i\}$ and
$\{g_i,h'_{ij},A'_i,\Phi'_i\}$ with cohomologous cocycles,
$h'_{ij}=\lambda_i^{-1}h_{ij}\lambda_j$ for smooth $\lambda_i\colon U_i\to G$, together with
$A'_i=\Ad_{\lambda_i^{-1}}A_i+\lambda_i^*\theta_G$ and $\Phi'_i=\rho(\lambda_i^{-1})\Phi_i$. Then
the maps $\varphi_i'^{-1}(x,a)\mapsto\varphi_i^{-1}(x,\lambda_i(x)a)$ agree on overlaps---exactly
because the cocycles are cohomologous---and assemble into a $G$-equivariant bundle isomorphism
$P'\to P$ covering $\id_M$, carrying $A'$ to $A$ and $\Phi'$ to $\Phi$. That is, it is a vertical
arrow of $\Ggeo$. Conversely, an arrow of $\Ggeo$ between global tuples, read in local
trivializations, exhibits the two systems of local data as related in exactly this way. Hence the
global tuple is well defined precisely as an object of $\Ggeo$ up to isomorphism.

\emph{Converse.}
Given a global tuple, choose any cover $\{U_i\}$ trivializing $P$ (possible since $M$ is
paracompact and $P$ is locally trivial); the induced $g_i:=g|_{U_i}$, $h_{ij}$, $A_i=s_i^*A$,
$\Phi_i$ satisfy \eqref{eq:cocycle}--\eqref{eq:matter-overlap} by the computations above.
\end{proof}

\begin{remark}[Changing the coordinate atlas is not a base diffeomorphism]
\label{rem:atlas}
The uniqueness clause above is stated for changes of \emph{local trivialization}, which induce
\emph{vertical} bundle isomorphisms covering $\id_M$. A change of \emph{coordinate atlas} is a
a different operation: re-charting $M$ changes the component arrays
$\check g_i$ of \eqref{eq:components} according to \eqref{eq:component-law} while leaving the
intrinsic tensor $g$---and hence the object $\Psi$ and its class $[\Psi]\in\Sgeo$---entirely
unchanged. It does not automatically produce a nontrivial base diffeomorphism, and it does not
move a point of $\Sgeo$. Nontrivial base diffeomorphisms enter only through the arrows of $\Ggeo$
with $f\neq\id_M$, which act on the geometry itself and not merely on its coordinate presentation.
\end{remark}

\begin{remark}[Descent and time orientation]
\label{rem:descent-time}
Theorem~\ref{thm:descent} glues the local metrics to a global Lorentzian $g$, but it does not
produce the time orientation $\mathfrak{t}$ that Definition~\ref{def:admissible} demands. Whether
$g$ admits any time orientation at all is a global topological condition on
$(M,g)$ (it holds if and only if the timelike cone bundle of $g$ has a global continuous section
of one of its components), and no purely local compatibility condition can force it. Two
consistent options are available, and a model must declare one.

\emph{Option (a): impose it after descent.} Glue $g$ as above, then require, as an additional
global admissibility condition, that $g$ be time-orientable and that a time orientation
$\mathfrak{t}$ be chosen and recorded in $\mathcal{B}$. Under this option
Theorem~\ref{thm:descent} reconstructs the metric, and membership of the reconstructed
configuration in $\Cgeo$ is a further global condition, verified separately.

\emph{Option (b): descend it.} Adjoin to the local data a choice, for each $i$, of a
$g_i$-timelike vector field $\tau_i\in\Gamma(TU_i)$, and impose the additional compatibility
condition that on every nonempty overlap $\tau_i$ and $\tau_j$ lie in the \emph{same} component
of the timelike cone,
\begin{equation}
\label{eq:time-glue}
g_i(\tau_i,\tau_j)<0 \qquad\text{on } U_i\cap U_j
\end{equation}
(with signature $(-,+,+,+)$, so that two timelike vectors have negative inner product exactly when
they are causally co-oriented). Condition \eqref{eq:time-glue} is a genuine gluing condition: the
relation ``lies in the same cone component as'' is an equivalence relation on timelike vectors at
each point, and the local choices $\tau_i$ define, by \eqref{eq:time-glue}, a coherent global
selection of one component of the cone bundle at every $x\in M$. That selection is by definition a
time orientation $\mathfrak{t}$ of the glued $g$, and it is unique given the $\tau_i$. Descent then
reconstructs the pair $(g,\mathfrak{t})$, and no further global condition on time orientation is
needed. (A partition of unity subordinate to $\{U_i\}$ shows that $\mathfrak{t}$ is realized by an
actual global timelike vector field, since the cone components are convex.)
\end{remark}

\begin{corollary}[Realizability boundary]
\label{cor:realizability-boundary}
A collection of proposed local metric, connection, and matter expressions represents a global
geometric source configuration only if the cocycle and overlap conditions
\eqref{eq:cocycle}--\eqref{eq:matter-overlap} hold on all overlaps, and only if the resulting
global tuple lies in the declared admissibility class $\Cgeo$---which, by
Remark~\ref{rem:descent-time}, includes the global time-orientation condition. Consequently, an
abstract tuple of relation values that admits no compatible descent data lies outside the
geometric realization domain: it is a point of the abstract product
$T\times\prod_\alpha V_\alpha$ that is not in the image of any geometric evaluation map.
\end{corollary}

\begin{proof}
Immediate from Theorem~\ref{thm:descent}: any global configuration restricts to compatible local
data, so failure of compatibility for \emph{every} presentation of the proposed values is an
obstruction to the existence of a configuration presenting them. Membership in $\Cgeo$ is a
further restriction, since Theorem~\ref{thm:descent} produces a global tuple but does not certify
that it satisfies the declared causal, boundary, or topological conditions.
\end{proof}

Corollary~\ref{cor:realizability-boundary} is the first and most basic restriction on the abstract
completion of Paper~3: \emph{its coordinates must be jointly realizable by global geometric
objects.} Two further restrictions---invariance (Section~\ref{sec:invariants}) and separation
(Section~\ref{sec:realization})---follow.

\begin{remark}[Obstructions are genuine]
\label{rem:obstruction-genuine}
That the constraint is not vacuous is visible already in cohomological terms: the isomorphism
class of $P$ is the class of $\{h_{ij}\}$ in the pointed set $\check H^1(M,\underline{G})$
(Remark~\ref{rem:cech}), and characteristic numbers computed from $A$ are constrained by that
class \cite{ChernWeil,MilnorStasheff}. A demand such as ``a $U(1)$ bundle over $S^2$ with first
Chern number $\tfrac12$'' fails at the cocycle stage; there is no compatible descent datum, hence
no configuration, hence no point of $\Sgeo$, hence no point of the realized completion. The
Cartesian product of independently nominated relation values is, in general, strictly larger than
what geometry attains.
\end{remark}

% ======================================================================
\section{Curvature, holonomy, and invariant global relations}
\label{sec:invariants}

\subsection{Curvature and the Bianchi identity}

The curvature of the connection $A$ is the $\gfrak$-valued $2$-form
\begin{equation}
\label{eq:curv}
F_A \;=\; \mathrm{d}A + A\wedge A
\qquad
\bigl(\text{equivalently } F_A=\mathrm{d}A+\tfrac12[A\wedge A]\bigr),
\end{equation}
in the matrix convention of Definition~\ref{def:compatible}. Under a gauge transformation
represented locally by $h\colon U\to G$,
\begin{equation}
\label{eq:gauge-transf}
A \;\longmapsto\; A^{h}=h^{-1}Ah+h^{-1}\mathrm{d}h,
\qquad
F_A \;\longmapsto\; F_{A^h}=h^{-1}F_A\,h ,
\end{equation}
so $F_A$ transforms in the adjoint representation: it is a section of
$\Lambda^2T^*M\otimes \mathrm{ad}P$, and is \emph{not} itself gauge invariant, though its
$\Ad$-invariant polynomial functions are. The Bianchi identity is
\begin{equation}
\label{eq:bianchi}
D_A F_A \;=\; \mathrm{d}F_A + [A,F_A] \;=\;0 ,
\end{equation}
an identity, not an equation of motion \cite{KobayashiNomizu1}.

The Levi-Civita connection $\nabla$ of $g$ is the unique affine connection with
\begin{equation}
\label{eq:LC}
\nabla g=0,
\qquad
\Tor(\nabla)=0 ,
\end{equation}
with Riemann tensor $R^{\rho}{}_{\sigma\mu\nu}$, Ricci tensor $R_{\mu\nu}=R^{\rho}{}_{\mu\rho\nu}$,
scalar curvature $R=g^{\mu\nu}R_{\mu\nu}$, and Einstein tensor
$G_{\mu\nu}=R_{\mu\nu}-\tfrac12 R\,g_{\mu\nu}$ \cite{Wald}.

\subsection{Invariant global relations}

\begin{definition}[Invariant geometric relation]
\label{def:invariant-relation}
A map $J\colon\Cgeo\to V$ into a set $V$ is an \emph{invariant geometric relation} if it is
constant on the arrows of $\Ggeo$, i.e.\ $J(\Psi)=J(\Psi')$ whenever there is an arrow
$\Psi\to\Psi'$. Equivalently, $J$ factors uniquely through $q\colon\Cgeo\to\Sgeo$, defining
$\bar J\colon\Sgeo\to V$. We write $J$ for both maps when no confusion arises.
\end{definition}

Invariance under $\Ggeo$ means invariance under \emph{both} gauge transformations and allowed
diffeomorphisms (including preservation of orientation and time orientation). Both are required;
neither alone suffices.

\begin{example}[Integral invariants]
\label{ex:integral-invariants}
Assume $M$ compact without boundary, or that $\Cgeo$ imposes decay/boundary conditions rendering
the following integrals absolutely convergent. Then each of
\begin{equation}
\label{eq:integral-invariants}
J_{\mathrm{EH}}(\Psi)=\int_M R_g\,\dvol,
\qquad
J_{\mathrm{YM}}(\Psi)=\int_M \bigl\langle F_A\wedge \ast F_A\bigr\rangle,
\qquad
J_{\mathrm{top}}(\Psi)=\int_M \bigl\langle F_A\wedge F_A\bigr\rangle
\end{equation}
is an invariant geometric relation. Gauge invariance of the latter two follows from
\eqref{eq:gauge-transf} with the $\Ad$-invariance of $\langle\cdot,\cdot\rangle$; diffeomorphism
invariance follows because each integrand is a top-degree form built naturally from the fields and
$\int_M f^*\omega=\int_M\omega$ for orientation-preserving $f$ (the Hodge star is natural:
$f^*(\ast_{g'}\eta)=\ast_{g}(f^*\eta)$ when $f_*g=g'$ and $f$ preserves orientation).
Concerning normalization, we make no claim beyond the following, which is all that is used: for
$M$ compact oriented without boundary and $G$ compact, the de Rham class of the Chern--Weil form
$\langle F_A\wedge F_A\rangle$ is independent of $A$ and represents (up to a normalization
constant depending on the chosen $\Ad$-invariant form $\langle\cdot,\cdot\rangle$ and on the
conventions for the Chern--Weil homomorphism) a characteristic class of $P$ in
$H^4_{\mathrm{dR}}(M;\Real)$ \cite{ChernWeil,MilnorStasheff,KobayashiNomizu2}. Hence
$J_{\mathrm{top}}$ is constant on each topological sector. We do \emph{not} assert a specific
integrality normalization here; any such normalization is a convention that must be fixed
together with $\langle\cdot,\cdot\rangle$, and none of the results below depend on the choice.
\end{example}

\begin{example}[Spectral relations]
\label{ex:spectral}
Let $\mathcal{D}[\Psi]$ be an operator constructed naturally from the configuration data
$(g,A,\Phi,\mathcal{B})$, together with a \emph{specified domain} determined by $\mathcal{B}$.
Then any function of
$\mathcal{D}[\Psi]$ that is invariant under isomorphism is an invariant geometric relation. The
spectrum, when it is well defined, is the leading example. The precise statement requires care,
and we distinguish two cases.

\emph{(i) Intertwining.} An arrow $F\colon\Psi\to\Psi'$ of $\Ggeo$ induces a linear isomorphism
$\mathcal{U}_F$ from the space of sections over $\Psi$ to the space of sections over $\Psi'$,
carrying the domain of $\mathcal{D}[\Psi]$ onto the domain of $\mathcal{D}[\Psi']$ and satisfying
the \emph{operator intertwining} relation
\begin{equation}
\label{eq:intertwine}
\mathcal{D}[\Psi']\circ\mathcal{U}_F=\mathcal{U}_F\circ\mathcal{D}[\Psi] .
\end{equation}
This holds by naturality of the construction and requires no metric structure on $W$. It already
implies that any relation depending only on the operator up to linear isomorphism (its spectrum as
a set, when the operator is, say, closed with nonempty resolvent set) is $\Ggeo$-invariant.

\emph{(ii) Unitary equivalence.} If, in addition, the representation $\rho$ is unitary (or, more
generally, $W$ carries a $\rho$-invariant Hermitian inner product), the associated bundle carries
a bundle metric preserved by $\hat F_\rho$, the base measure $\dvol$ is preserved by the allowed
$f$, and the declared boundary conditions define a self-adjoint realization of $\mathcal{D}$ on the
resulting Hilbert space, then $\mathcal{U}_F$ is \emph{unitary} and \eqref{eq:intertwine} is a
unitary equivalence. In that case the spectrum as a \emph{multiset} (with multiplicities), the
heat trace $\tr e^{-t\mathcal{D}}$, zeta values, and eigenvalue gaps are all invariant geometric
relations.

An arbitrary finite-dimensional $\rho\colon G\to\mathrm{GL}(W)$ is not unitary, so case (ii)
requires its stated hypotheses; absent them, only the intertwining statement (i) is available. In
either case the specification of the domain is essential: without it the spectrum is not a function
of $[\Psi]$, and distinct self-adjoint extensions yield distinct invariant relations.
\end{example}

\subsection{Holonomy and Wilson observables}

Let $\gamma$ be a piecewise-smooth closed curve in $M$ based at $x\in M$, and let $p\in P_x$.
Parallel transport of $A$ along $\gamma$ starting at $p$ ends at $p\cdot a$ for a unique $a\in G$;
writing $\Hol_\gamma^{p}(A):=a$, a change of starting point $p\mapsto p\cdot b$ gives
$\Hol_\gamma^{p\cdot b}(A)=b^{-1}\Hol_\gamma^p(A)\,b$. In a local trivialization with gauge $h$,
this reads
\begin{equation}
\label{eq:hol-conj}
\Hol_\gamma(A^{h})\;=\;h(x)^{-1}\,\Hol_\gamma(A)\,h(x) .
\end{equation}

\begin{proposition}[Gauge invariance of holonomy class and Wilson observables]
\label{prop:wilson-gauge}
For a fixed based loop $\gamma$, the conjugacy class
\begin{equation}
\label{eq:hol-class}
\bigl[\Hol_\gamma(A)\bigr]_{\mathrm{conj}} \in G/\!\!\sim_{\mathrm{conj}}
\end{equation}
is independent of the choice of $p\in P_x$ and is invariant under all gauge transformations.
Consequently, for any finite-dimensional representation $\sigma$ of $G$, the Wilson observable
\begin{equation}
\label{eq:wilson}
W_{\gamma,\sigma}(A)\;=\;\tr_\sigma \Hol_\gamma(A)
\end{equation}
is gauge invariant. If $A$ is flat, $W_{\gamma,\sigma}$ depends on $\gamma$ only through its free
homotopy class.
\end{proposition}

\begin{proof}
Independence of $p$ and gauge invariance are immediate from \eqref{eq:hol-conj}, since a conjugacy
class is invariant under conjugation and $\tr_\sigma$ is a class function. The final statement
holds because for flat $A$ parallel transport depends only on the homotopy class of the path rel
endpoints, and a free homotopy of based loops changes the holonomy by conjugation, which
$\tr_\sigma$ does not see \cite{KobayashiNomizu1}.
\end{proof}

\subsection{Loops and diffeomorphism invariance}
\label{sec:diff-loops}

Proposition~\ref{prop:wilson-gauge} concerns \emph{gauge} invariance only. A Wilson observable
attached to a \emph{fixed} loop $\gamma\subset M$ is in general \emph{not} invariant under the
allowed diffeomorphisms, because a diffeomorphism $f$ moves the loop:
$W_{\gamma,\sigma}(f_*A)=W_{f^{-1}\circ\gamma,\sigma}(A)$, and $f^{-1}\circ\gamma\neq\gamma$ in
general. Three constructions restore genuine $\Ggeo$-invariance. They are mutually exclusive as
declarations, and a model must state which it adopts; each is used somewhere below.

\medskip
\noindent\textbf{(L1) Include the loop in the configuration.}
Enlarge the tuple \eqref{eq:config} to $(M,g,P,A,\Phi,\mathcal{B},\gamma)$ and extend the arrows
of $\Ggeo$ to carry $\gamma$ to $f\circ\gamma$. Then
$\Psi\mapsto W_{\gamma,\sigma}(A)$ is invariant on the enlarged groupoid by construction:
$W_{f\circ\gamma,\sigma}(f_*A)=W_{\gamma,\sigma}(A)$, since parallel transport is natural under
bundle isomorphisms. The loop is then part of the source data, not a background structure, and
$\Sgeo$ is correspondingly enlarged.

\medskip
\noindent\textbf{(L2) Use an invariant constructed from the whole loop space.}
Examples: the \emph{set} of Wilson values $\{W_{\gamma,\sigma}(A):\gamma\in\Omega(M,x)\}\subseteq\Cx$
(a subset of $\Cx$, hence unaffected by any relabelling of loops); the holonomy group
$\mathrm{Hol}_x(A)\le G$ up to conjugacy; or, for flat $A$, the holonomy homomorphism
$\pi_1(M,x)\to G$ taken modulo conjugation \emph{and} modulo the outer action of $f_*$ on
$\pi_1(M,x)$. Each is $\Ggeo$-invariant without enlarging the configuration, precisely because the
diffeomorphism permutes the loops over which the construction ranges.

\medskip
\noindent\textbf{(L3) Let the declared data pin a homology class.}
If $\mathcal{B}$ distinguishes a class $c\in H_1(M;\Int)$ that every allowed diffeomorphism
preserves, and if the observable depends on $\gamma$ only through $c$, then it is invariant. This
is what happens in Section~\ref{sec:u1}: there $G=U(1)$ is abelian and $A$ is flat, so
$\Hol_\gamma(A)$ depends on $\gamma$ only through its class in $H_1(S^1;\Int)\cong\Int$; the
allowed diffeomorphisms are declared to be the orientation-preserving ones, which act trivially on
$H_1(S^1;\Int)$; hence $\Hol_{\gamma_0}(A)$ is a genuine function on $\Sgeo$. (Under an
orientation-\emph{reversing} diffeomorphism the generator is sent to its inverse and
$\Hol\mapsto\Hol^{-1}$; this is exactly why the orientation is part of the declared data.)

\begin{remark}[Curvature does not determine transport]
\label{rem:curv-vs-hol}
Local curvature and global transport are related but not interchangeable. Curvature is the
infinitesimal holonomy: the Ambrose--Singer theorem \cite{AmbroseSinger,KobayashiNomizu1} states
that the Lie algebra of the holonomy group of a connection on a connected principal bundle is
spanned by the curvature values $F_A(X,Y)$ at points reachable from the base point by horizontal
curves, transported back along those curves. Nevertheless holonomy around noncontractible cycles is
not determined by $F_A$: flat connections with $F_A\equiv 0$ can have arbitrary holonomy
homomorphisms $\pi_1(M)\to G$ up to conjugacy. This is not a formal curiosity; it is the content of
the Aharonov--Bohm effect \cite{AharonovBohm} and of the Wu--Yang analysis of the global
significance of gauge potentials \cite{WuYang}, in which the physically operative datum is a
holonomy in a region where the field strength vanishes identically. Section~\ref{sec:u1} is the
sharpest mathematical illustration: there $F_A\equiv0$ identically, and yet the reduced source is a
full circle of physically inequivalent configurations, completely parametrized by holonomy.
\end{remark}

\begin{remark}[No finite separating list of curvature scalars]
\label{rem:no-finite-separation}
It is not assumed anywhere below that a finite family of curvature scalars---or any fixed finite
family of local invariants---separates arbitrary Lorentzian geometries. It does not: in Lorentzian
signature there exist non-flat metrics all of whose polynomial scalar curvature invariants vanish
identically, and which therefore cannot be distinguished from flat space by any such family; the
plane-wave and more general Kundt-type ``VSI'' spacetimes are the standard examples
\cite{Pravda2002}. Separation is therefore never invoked as a general fact. Where it is needed it
is proved on the restricted domain of the model, as in Theorem~\ref{thm:hol-class}.
\end{remark}

% ======================================================================
\section{Geometric realization of the minimal completion}
\label{sec:realization}

\subsection{The geometric evaluation map}

\begin{definition}[Declared readout and relation family]
\label{def:readout}
Let
\begin{equation}
\label{eq:pgeo}
\pgeo\colon \Sgeo\longrightarrow T
\end{equation}
be a declared Tier-1 readout: a map, defined on the reduced source, recording exactly the data
retained by the Tier-1 description of the model. Let $A$ be an index \emph{set} and let
\begin{equation}
\label{eq:Jalpha}
J_\alpha\colon \Sgeo\longrightarrow V_\alpha,
\qquad \alpha\in A,
\end{equation}
be invariant geometric relations in the sense of Definition~\ref{def:invariant-relation}. Write
$J_A=(J_\alpha)_{\alpha\in A}$ and
\begin{equation}
\label{eq:jgeo}
\jgeo \;=\;\bigl(\pgeo,\;J_A\bigr)\colon
\Sgeo\longrightarrow T\times\prod_{\alpha\in A}V_\alpha .
\end{equation}
\end{definition}

\begin{definition}[Geometrically realized completion]
\label{def:egeo}
The \emph{geometrically realized completion} of $\pgeo$ relative to $\{J_\alpha\}_{\alpha\in A}$ is
\begin{equation}
\label{eq:egeo}
\Egeo \;:=\; \Imm \jgeo \;\subseteq\; T\times\prod_{\alpha\in A}V_\alpha ,
\end{equation}
with corestriction $\egeo\colon\Sgeo\twoheadrightarrow \Egeo$ and forgetful map
\begin{equation}
\label{eq:uA}
u_A\colon \Egeo\longrightarrow T,
\qquad
u_A\bigl(t,(v_\alpha)\bigr)=t ,
\end{equation}
so that $u_A\circ \egeo=\pgeo$ by construction.
\end{definition}

Note that $\Egeo$ is the image, \emph{not} the ambient product. This restriction performs the
essential mathematical work: it enforces all the compatibility relations among metric, bundle,
connection, matter, boundary, and invariant data that Section~\ref{sec:descent} exhibited. The
Cartesian product of independently written coordinate values is in general \emph{not} realizable
(Corollary~\ref{cor:realizability-boundary}, Remark~\ref{rem:obstruction-genuine}).

\subsection{The restricted realization theorem}

\begin{theorem}[Restricted geometric realization]
\label{thm:restricted-realization}
Fix an admissibility class $\Cgeo$, a readout $\pgeo$, and a family $\{J_\alpha\}_{\alpha\in A}$ of
invariant geometric relations on $\Sgeo$. Then:
\begin{enumerate}[label=(\arabic*),leftmargin=2.2em]
\item \emph{(Realizability.)} For every $\xi\in\Egeo$ there exists $[\Psi]\in\Sgeo$ with
      $\egeo([\Psi])=\xi$.
\item \emph{(Factorization.)} $\pgeo=u_A\circ\egeo$ and
      $J_\alpha=\mathrm{pr}_\alpha\circ\egeo$ for every $\alpha$.
\item \emph{(Terminality.)} If $\pi\colon\Sgeo\twoheadrightarrow E$ is any surjection through which
      $\pgeo$ and every $J_\alpha$ factor---say $\pgeo=u\circ\pi$ and
      $J_\alpha=\kappa_\alpha\circ\pi$---then there is a \emph{unique} map
      $\theta\colon E\to\Egeo$ with $\theta\circ\pi=\egeo$; $\theta$ is surjective and satisfies
      $u_A\circ\theta=u$ and $\mathrm{pr}_\alpha\circ\theta=\kappa_\alpha$.
\end{enumerate}
Thus $\Egeo$ is the canonical minimal completion of $\pgeo$ relative to $\{J_\alpha\}$, in the
sense of Paper~3 \cite{ST3}, computed on the geometric source $\Sgeo$.
\end{theorem}

\begin{proof}
(1) is immediate: $\Egeo$ is by definition the image of $\jgeo$. (2) holds by definition of $\jgeo$
and of the projections; well-definedness of $\pgeo$ and each $J_\alpha$ \emph{as maps on $\Sgeo$}
is exactly the invariance hypothesis of Definition~\ref{def:invariant-relation}, which is where
gauge and diffeomorphism invariance are consumed.

(3) Define $\theta(e):=\egeo(s)$ for any $s\in\pi^{-1}(e)$; the fibre is nonempty since $\pi$ is
surjective. This is well defined: if $\pi(s)=\pi(s')=e$ then
$\pgeo(s)=u(\pi(s))=u(\pi(s'))=\pgeo(s')$ and likewise
$J_\alpha(s)=\kappa_\alpha(\pi(s))=\kappa_\alpha(\pi(s'))=J_\alpha(s')$ for every $\alpha$; hence
$\jgeo(s)=\jgeo(s')$, i.e.\ $\egeo(s)=\egeo(s')$. By construction $\theta\circ\pi=\egeo$, and since
$\egeo$ is surjective onto $\Egeo$, so is $\theta$. Uniqueness: any $\theta'$ with
$\theta'\circ\pi=\egeo$ agrees with $\theta$ on $\pi(\Sgeo)=E$, hence everywhere. The
compatibilities follow by evaluating on $\pi(s)$ and using surjectivity of $\pi$.

This is the terminal universal property established in Paper~3 \cite{ST3}, applied to the source set
$\Sgeo$, the readout $\pgeo$, and the relation family $\{J_\alpha\}$; the additional content here is
that $\Sgeo$ is the \emph{reduced} geometric source and that the $J_\alpha$ are \emph{invariant}
geometric relations, so that all maps involved are functions of physical orbits rather than of gauge
or coordinate representatives.
\end{proof}

\begin{remark}[Relative, not absolute]
\label{rem:relative}
Theorem~\ref{thm:restricted-realization} is \emph{restricted} in two independent senses, both
essential. First, it realizes the completion only \emph{relative to the declared family}
$\{J_\alpha\}$: a different family gives a different $\Egeo$, generally neither finer nor coarser.
Second, it realizes only what geometry attains: $\Egeo$ is an image, and
Corollary~\ref{cor:realizability-boundary} shows that this image can be a proper subset of the
abstract product. Neither the theorem nor anything else in this paper asserts that an arbitrarily
nominated abstract completion $E_A$ of Paper~3 is geometrically realizable. It is realizable exactly
when it is isomorphic, over $T$, to $\Egeo$ for some admissible $\Cgeo$, readout, and invariant
relation family---and establishing that in a given model is work.
\end{remark}

\subsection{When does the completion reconstruct the source?}

\begin{proposition}[Orbit reconstruction criterion]
\label{prop:orbit-reconstruction}
The completion map $\egeo\colon\Sgeo\to\Egeo$ is injective---equivalently, bijective onto
$\Egeo$---if and only if the family
\begin{equation}
\label{eq:sep-family}
\{\pgeo\}\cup\{J_\alpha : \alpha\in A\}
\end{equation}
\emph{separates points of $\Sgeo$}: for all $[\Psi],[\Psi']\in\Sgeo$, equality of the readouts and
of every $J_\alpha$ forces $[\Psi]=[\Psi']$. In that case $\Sgeo\cong\Egeo$ as sets over $T$.
\end{proposition}

\begin{proof}
By construction $\egeo([\Psi])=\egeo([\Psi'])$ if and only if the readouts agree and every
$J_\alpha$ agrees. Injectivity of $\egeo$ is therefore literally the separation property.
Surjectivity onto $\Egeo$ always holds, by Theorem~\ref{thm:restricted-realization}(1). The
identification is over $T$ because $u_A\circ\egeo=\pgeo$.
\end{proof}

Proposition~\ref{prop:orbit-reconstruction} converts a question about completions into a question of
invariant theory: \emph{does the declared family separate the orbits?} No general answer exists, and
none is asserted (Remark~\ref{rem:no-finite-separation}). Separation must be proved on the restricted
domain relevant to the model. Section~\ref{sec:u1} carries out exactly such a proof in a case where
the answer is affirmative and complete.

% ======================================================================
\section{Classification of flat \texorpdfstring{$U(1)$}{U(1)} connections on an oriented circle}
\label{sec:u1}

This section supplies an exact classification example: a restricted domain on which the separation
hypothesis of Proposition~\ref{prop:orbit-reconstruction} can be \emph{proved}, and on which the
geometrically realized completion can be identified on the nose.

\subsection{The model and its conventions}

\begin{convention}[Scope of the circle model]
\label{conv:circle-scope}
This section is a deliberate one-dimensional specialization, and two stipulations of the general
framework are suspended for it. Both are recorded here, since leaving either implicit would create a
genuine inconsistency with Section~\ref{sec:source}.

\emph{Dimension.} The base is $M=S^1_L$, so $\dim M=1$ rather than the $\dim M=4$ of
Convention~\ref{conv:basic}, and the metric is Riemannian rather than Lorentzian. Nothing in
Sections~\ref{sec:source}--\ref{sec:realization} used the dimension or the signature: the groupoid,
the orbit space, Definition~\ref{def:invariant-relation}, Theorem~\ref{thm:restricted-realization},
and Proposition~\ref{prop:orbit-reconstruction} are all dimension- and signature-agnostic, and it is
only these that the present section invokes. The descent theorem and the variational apparatus of
Sections~\ref{sec:variational}--\ref{sec:noether} are not used here.

\emph{The metric and the allowed maps.} A tension must be resolved explicitly. The model fixes the
standard metric $\mathrm{d}x^{2}$ on $S^1_L$ while also admitting \emph{every}
orientation-preserving diffeomorphism as an allowed base map; under the groupoid of
Definition~\ref{def:groupoid}, whose arrows satisfy $f_*g=g'$, a fixed metric would instead restrict
the allowed base maps to its isometries, namely the rigid rotations. Of the three consistent ways to
proceed---restricting the allowed diffeomorphisms to orientation-preserving isometries; transforming
the metric along with the connection; or defining the circle source purely as the gauge quotient and
treating the diffeomorphism question separately---we adopt the third, which is the formulation the
arguments below actually use.

Accordingly, in this section the metric $\mathrm{d}x^{2}$ and its measure $\mathrm{d}x$ are treated
as \emph{fixed background structure}, not as part of the configuration: the configuration is the
connection $A$ alone, and the reduced circle source is \emph{defined} to be the gauge quotient
\begin{equation}
\label{eq:circle-source-def}
\Sgeo \;:=\; \Cgeo^{\,\circ}\big/\Gau(P) .
\end{equation}
The background metric is what gives the operator $\mathcal{D}_A$ of \eqref{eq:Dop} its meaning, and
it is not transformed by the maps considered below. The role of the orientation-preserving
diffeomorphisms is then a separate matter, treated separately: they are \emph{not} arrows acting on
a metric they would have to preserve, but are used solely through the fact that they fix the
homology generator $c$, and hence preserve holonomy. Lemma~\ref{lem:spec-invariance}(1) shows that
they consequently identify no classes beyond gauge equivalence, so that \eqref{eq:circle-source-def}
loses nothing by omitting them; and Lemma~\ref{lem:spec-invariance}(2) shows that the spectral
relation is invariant under them for the same reason. Definition~\ref{def:circle-model} is to be
read in this light.
\end{convention}

\begin{definition}[The circle model]
\label{def:circle-model}
Fix $L>0$ and let $M=S^1_L:=\Real/L\Int$ with its standard orientation and the global angular
coordinate $x$ (well defined modulo $L$). Let $G=U(1)$ and let $P=S^1_L\times U(1)$ be the trivial
bundle; this is no loss of generality, since every principal $U(1)$-bundle over $S^1$ is trivial
(the classifying space argument gives
$\bigl[S^1,BU(1)\bigr]\cong H^2(S^1;\Int)=0$; equivalently, every $U(1)$-bundle over a
one-complex is trivial \cite{Steenrod}). Let $\Cgeo^{\,\circ}$ be the class of smooth connections
$A$ on $P$, and let $\Gau(P)=C^\infty(S^1_L,U(1))$ be the gauge group; by
Convention~\ref{conv:circle-scope} the configuration is the connection alone, and the reduced source
is the gauge quotient \eqref{eq:circle-source-def}. The declared data $\mathcal{B}$ consist of the
orientation of $S^1_L$, the resulting generator
$c=[\gamma_0]\in H_1(S^1_L;\Int)\cong\Int$ of the fundamental cycle, $\gamma_0(s)=s$, $s\in[0,L]$,
and the fixed background metric $\mathrm{d}x^{2}$ with its measure. The orientation-preserving
diffeomorphisms of $S^1_L$ fix $c$; holonomy around $c$ is therefore a legitimate function on orbits
by mechanism (L3) of Section~\ref{sec:diff-loops}, and its invariance under those diffeomorphisms is
established in Lemma~\ref{lem:spec-invariance}.
\end{definition}

Write a connection in the global trivialization as
\begin{equation}
\label{eq:A-circle}
A=a(x)\,\mathrm{d}x,
\qquad a\in C^\infty(S^1_L,\Real) .
\end{equation}
Because $\dim M=1$, every $2$-form on $M$ vanishes; hence
\begin{equation}
\label{eq:flat}
F_A=\mathrm{d}A\equiv 0 \qquad\text{for every }A .
\end{equation}
\emph{Every} connection on the circle is flat. Local curvature is identically zero on this entire
domain and therefore carries no information whatsoever; whatever separates orbits here must be
global.

\begin{convention}[Gauge action and charged covariant derivative]
\label{conv:u1-norm}
Identify $\mathfrak{u}(1)=i\Real$ and write connections with the real potential $a$ as in
\eqref{eq:A-circle}. A gauge transformation is a smooth map $u\colon S^1_L\to U(1)$, acting on the
connection by
\begin{equation}
\label{eq:u1-gauge}
A\;\longmapsto\;A^{u}\;=\;A-i\,u^{-1}\mathrm{d}u ,
\end{equation}
which is real-valued because $-i\,u^{-1}\mathrm{d}u$ is real for $|u|=1$. On a charged scalar
$\phi\in C^\infty(S^1_L,\Cx)$ of unit charge, the same $u$ acts by
\begin{equation}
\label{eq:u1-matter}
\phi\;\longmapsto\;\phi^{u}\;=\;u\,\phi ,
\end{equation}
and the charged covariant derivative is
\begin{equation}
\label{eq:u1-covder}
D_A\phi\;=\;\bigl(\partial_x - i\,a(x)\bigr)\phi\;\mathrm{d}x .
\end{equation}
These three conventions are mutually consistent. Writing $u=e^{i\varphi}$ locally, so that
$-i\,u^{-1}\mathrm{d}u=\mathrm{d}\varphi$ and hence $a^{u}=a+\varphi'$, we get
\[
\bigl(\partial_x-i\,a^{u}\bigr)\phi^{u}
=\bigl(\partial_x - i(a+\varphi')\bigr)\bigl(e^{i\varphi}\phi\bigr)
= e^{i\varphi}\Bigl(i\varphi'\phi+\partial_x\phi - i a\phi - i\varphi'\phi\Bigr)
= u\,\bigl(\partial_x - ia\bigr)\phi ,
\]
so $D_{A^u}\phi^{u}=u\,D_A\phi$ transforms covariantly, as required.
\end{convention}

\begin{definition}[Holonomy on the circle]
\label{def:hol-circle}
For $A=a\,\mathrm{d}x$ set
\begin{equation}
\label{eq:hol-circle}
\Hol(A)\;:=\;\exp\!\left(i\int_0^L a(x)\,\mathrm{d}x\right)\;\in\;U(1),
\end{equation}
the holonomy around the declared fundamental cycle $\gamma_0$.
\end{definition}

\subsection{The classification theorem}

\begin{theorem}[Holonomy classification]
\label{thm:hol-class}
Two smooth $U(1)$ connections $A=a\,\mathrm{d}x$ and $A'=a'\,\mathrm{d}x$ on $S^1_L$ are gauge
equivalent if and only if
\begin{equation}
\label{eq:hol-equal}
\Hol(A)=\Hol(A') .
\end{equation}
Consequently $[A]\mapsto\Hol(A)$ is a bijection
\begin{equation}
\label{eq:moduli}
\Mflat\bigl(S^1,U(1)\bigr)
\;:=\;
\Cgeo^{\,\circ}/\Gau(P)
\;\;\xrightarrow{\ \cong\ }\;\;
U(1),
\end{equation}
and it is a homeomorphism when the left-hand side carries the quotient of the $C^\infty$ topology.
\end{theorem}

\begin{proof}
\emph{Step 1: holonomy is gauge invariant.}
Let $u\in C^\infty(S^1_L,U(1))$. Pull $u$ back to the universal cover: the composite
$\tilde u:=u\circ\mathrm{pr}\colon[0,L]\to U(1)$, where $\mathrm{pr}(s)=s\bmod L$, lifts along the
covering $\Real\xrightarrow{\,t\mapsto e^{it}\,}U(1)$ to a smooth
$\varphi\colon[0,L]\to\Real$ with $u(x)=e^{i\varphi(x)}$ for $x\in[0,L]$, because $[0,L]$ is simply
connected. Single-valuedness of $u$ on $S^1_L$ forces
\begin{equation}
\label{eq:winding}
\varphi(L)-\varphi(0)\;=\;2\pi n,
\qquad n=\deg(u)\in\Int .
\end{equation}
Now $-i\,u^{-1}\mathrm{d}u = -i\,e^{-i\varphi}\cdot i\varphi'\,e^{i\varphi}\,\mathrm{d}x
=\varphi'\,\mathrm{d}x=\mathrm{d}\varphi$, so the declared action \eqref{eq:u1-gauge} gives
\begin{equation}
\label{eq:Au-explicit}
A^{u}\;=\;A-i\,u^{-1}\mathrm{d}u\;=\;A+\mathrm{d}\varphi,
\qquad\text{i.e.}\qquad
a^{u}=a+\varphi' .
\end{equation}
Hence the period changes by $+2\pi n$:
\[
\int_0^L a^{u}
=\int_0^L a+\bigl(\varphi(L)-\varphi(0)\bigr)
=\int_0^L a+2\pi n ,
\]
and therefore
\[
\Hol(A^{u})
=\exp\!\Bigl(i\int_0^L a + 2\pi i n\Bigr)
=\exp\!\Bigl(i\int_0^L a\Bigr)
=\Hol(A) .
\]
The gauge group can shift the period integral, but only by integer multiples of $2\pi$, which
$\exp(i\,\cdot\,)$ does not see. (The sign of the shift is $+2\pi n$ with the declared convention;
either sign would of course leave the holonomy unchanged, but we record the correct one.)

\emph{Step 2: equal holonomy implies gauge equivalence.}
Suppose $\Hol(A)=\Hol(A')$. Then $\exp\bigl(i\int_0^L(a'-a)\bigr)=1$, so there is $n\in\Int$ with
\begin{equation}
\label{eq:period-int}
\int_0^L\bigl(a'(x)-a(x)\bigr)\,\mathrm{d}x \;=\; 2\pi n .
\end{equation}
Define
\begin{equation}
\label{eq:chi}
\chi(x)\;:=\;\int_0^x\bigl(a'(y)-a(y)\bigr)\,\mathrm{d}y,
\qquad x\in[0,L],
\end{equation}
so that $\chi$ is smooth with
\begin{equation}
\label{eq:chi-prime}
\chi'=a'-a
\qquad\text{and}\qquad
\chi(L)-\chi(0)=2\pi n
\end{equation}
by \eqref{eq:period-int}. The function $\chi$ itself need \emph{not} descend to $S^1_L$, but
\begin{equation}
\label{eq:u-def}
u(x)\;:=\;e^{\,i\chi(x)}
\end{equation}
does: $u(L)=e^{i\chi(0)}e^{2\pi i n}=e^{i\chi(0)}=u(0)$, and all derivatives match at the seam
because $a,a'$ are smooth and $L$-periodic, so $\chi'$ and all its derivatives are $L$-periodic.
Hence $u\in C^\infty(S^1_L,U(1))$ is a bona fide gauge transformation, of degree $n$. By exactly
the computation of Step~1 applied to this $u$ (with $\varphi=\chi$),
\begin{equation}
\label{eq:u-mc}
-i\,u^{-1}\mathrm{d}u \;=\; \mathrm{d}\chi \;=\; (a'-a)\,\mathrm{d}x \;=\; A'-A ,
\end{equation}
using \eqref{eq:chi-prime}. Therefore, directly from the declared action \eqref{eq:u1-gauge} in the
form \eqref{eq:Au-explicit},
\[
A^{u}\;=\;A-i\,u^{-1}\mathrm{d}u\;=\;A+(A'-A)\;=\;A' .
\]
So $u$ itself carries $A$ to $A'$, and the two connections are gauge equivalent.

\emph{Step 3: the bijection.}
By Steps 1 and 2, $\Hol$ descends to an injection $\Cgeo^{\,\circ}/\Gau(P)\hookrightarrow U(1)$. It
is surjective: for $\vartheta\in\Real$ the constant connection
\begin{equation}
\label{eq:Atheta}
A_\vartheta:=\frac{\vartheta}{L}\,\mathrm{d}x
\qquad\text{satisfies}\qquad
\Hol(A_\vartheta)=e^{i\vartheta},
\end{equation}
and every element of $U(1)$ is of the form $e^{i\vartheta}$. Hence \eqref{eq:moduli} is a bijection.

\emph{Step 4: the homeomorphism.}
Give $\Cgeo^{\,\circ}\cong C^\infty(S^1_L,\Real)$ its usual Fr\'echet ($C^\infty$) topology and the
quotient $\mathcal{M}:=\Cgeo^{\,\circ}/\Gau(P)$ the quotient topology, with projection
$\mathrm{pr}_{\mathcal{M}}\colon\Cgeo^{\,\circ}\to\mathcal{M}$. Write
$\Theta\colon\mathcal{M}\to U(1)$ for the bijection \eqref{eq:moduli}.

$\Theta$ is continuous: the map $\Cgeo^{\,\circ}\to U(1)$, $A\mapsto\exp(i\int_0^L a)$, is
continuous (indeed it is continuous already in the $L^1$ topology, hence in $C^\infty$), and it is
constant on $\Gau(P)$-orbits by Step~1; the universal property of the quotient topology then makes
$\Theta$ continuous.

For the inverse, we do \emph{not} attempt to choose a global continuous argument $\vartheta$ on
$U(1)$---no such choice exists. Instead define
\begin{equation}
\label{eq:sigma}
\sigma\colon\Real\longrightarrow\mathcal{M},
\qquad
\sigma(\vartheta):=\mathrm{pr}_{\mathcal{M}}\bigl(A_\vartheta\bigr)
=\bigl[\tfrac{\vartheta}{L}\,\mathrm{d}x\bigr] .
\end{equation}
The map $\Real\to\Cgeo^{\,\circ}$, $\vartheta\mapsto A_\vartheta$, is continuous (it is linear into
the constants, and the $C^\infty$ seminorms of $A_\vartheta$ depend continuously on $\vartheta$), so
$\sigma$ is continuous as a composite with $\mathrm{pr}_{\mathcal{M}}$.

$\sigma$ is $2\pi$-periodic. Indeed $A_{\vartheta+2\pi}-A_\vartheta=\tfrac{2\pi}{L}\mathrm{d}x$, and
by Step~2 (with $a=\vartheta/L$, $a'=(\vartheta+2\pi)/L$, so that $\int_0^L(a'-a)=2\pi$, i.e.\
$n=1$) the connections $A_\vartheta$ and $A_{\vartheta+2\pi}$ are gauge equivalent; explicitly the
gauge transformation is $u(x)=e^{2\pi i x/L}$, of degree $1$. Hence
$\sigma(\vartheta+2\pi)=\sigma(\vartheta)$ for all $\vartheta$.

Therefore $\sigma$ factors through the quotient $\Real\to\Real/2\pi\Int$, and since
$\Real/2\pi\Int\cong U(1)$ via $\vartheta\bmod 2\pi\mapsto e^{i\vartheta}$ (a homeomorphism), we
obtain a continuous map
\begin{equation}
\label{eq:sigma-bar}
\bar\sigma\colon U(1)\longrightarrow\mathcal{M},
\qquad
\bar\sigma\bigl(e^{i\vartheta}\bigr)=\sigma(\vartheta) ,
\end{equation}
well defined and continuous by the universal property of the quotient $\Real\to\Real/2\pi\Int$
(which is an open, hence quotient, map). Finally $\Theta\circ\bar\sigma=\id_{U(1)}$ by
\eqref{eq:Atheta}, and $\bar\sigma\circ\Theta=\id_{\mathcal{M}}$ because both sides are the identity
on the classes $[A_\vartheta]$, which exhaust $\mathcal{M}$ by Step~3 together with the injectivity
of $\Theta$. Hence $\bar\sigma=\Theta^{-1}$ is continuous, and $\Theta$ is a homeomorphism.
\end{proof}

\begin{remark}[Where the hypotheses are used]
\label{rem:u1-hyps}
Three hypotheses were used and cannot be dropped. \emph{Smoothness and periodicity} of $a,a'$ give
smoothness of $u$ across the seam in Step~2. \emph{Connectedness of $S^1$}, with
$\pi_1(S^1)\cong\Int$ generated by a single class, is what makes one period integral a complete
invariant. \emph{Orientation}, fixed in Definition~\ref{def:circle-model}, is what makes $\int_0^L$
well defined with a sign: reversing orientation sends $\Hol\mapsto\Hol^{-1}$, and it is precisely
this ambiguity that reappears in Corollary~\ref{cor:spec}. The declaration of the allowed
diffeomorphism group as the orientation-preserving one is therefore part of the hypothesis of the
theorem.
\end{remark}

\subsection{The realized completions}

\begin{corollary}[Exact geometric realization of the holonomy completion]
\label{cor:hol-completion}
In the circle model, take the curvature readout
\begin{equation}
\label{eq:pF}
p_F\colon \Sgeo=\Mflat(S^1,U(1))\longrightarrow T=\{0\},
\qquad
p_F\bigl([A]\bigr)=F_A=0 ,
\end{equation}
which by \eqref{eq:flat} is constant and hence retains \emph{no} information, and take the single
invariant relation $J_{\mathrm{hol}}([A])=\Hol(A)\in U(1)$. Then
\begin{equation}
\label{eq:Ehol}
E_{\mathrm{hol}}
=\Imm\bigl(p_F,J_{\mathrm{hol}}\bigr)
\;\cong\;
\Mflat\bigl(S^1,U(1)\bigr)
\;\cong\; U(1),
\end{equation}
and the completion map $e_{\mathrm{hol}}\colon\Sgeo\to E_{\mathrm{hol}}$ is a \emph{bijection}.
\end{corollary}

\begin{proof}
$\Imm(p_F,J_{\mathrm{hol}})=\{0\}\times\Imm J_{\mathrm{hol}}=\{0\}\times U(1)\cong U(1)$, by the
surjectivity established in Step~3 of Theorem~\ref{thm:hol-class}. The family
$\{p_F\}\cup\{J_{\mathrm{hol}}\}$ separates orbits by the injectivity established there, so
Proposition~\ref{prop:orbit-reconstruction} gives injectivity of $e_{\mathrm{hol}}$; surjectivity is
Theorem~\ref{thm:restricted-realization}(1).
\end{proof}

Corollary~\ref{cor:hol-completion} is the sharpest illustration in this paper of the mechanism the
sequence is about. The readout $p_F$ is maximally degenerate---it is constant---so on its own it
distinguishes nothing. Yet the reduced source is not a point: it is a circle of genuinely
inequivalent physical configurations. The completion recovers exactly that circle, from a single
global relation, on a domain where local curvature vanishes identically. Holonomy here is not one
invariant among many: it separates \emph{all} gauge orbits and therefore reconstructs the reduced
geometric source completely. This is the mathematical content of the Aharonov--Bohm and Wu--Yang
observations \cite{AharonovBohm,WuYang}.

\begin{lemma}[The reduced circle source, and invariance of the charged spectrum]
\label{lem:spec-invariance}
In the circle model of Definition~\ref{def:circle-model}, with the scope fixed by
Convention~\ref{conv:circle-scope}:
\begin{enumerate}[label=(\arabic*),leftmargin=2.2em]
\item The orientation-preserving diffeomorphisms of $S^1_L$ identify no connection classes beyond
      gauge equivalence: for every such $f$ and every $A$, the pushforward $f_*A$ is gauge
      equivalent to $A$. Hence adjoining them to the gauge group would coarsen nothing, and the
      reduced source \eqref{eq:circle-source-def} may be written
      \begin{equation}
      \label{eq:Sgeo-circle}
      \Sgeo \;=\; \Cgeo^{\,\circ}\big/\Gau(P) \;=\; \Mflat\bigl(S^1,U(1)\bigr) .
      \end{equation}
\item The spectral multiset $\operatorname{spec}\mathcal{D}_A$ of \eqref{eq:Dop}, computed with
      respect to the fixed background metric $\mathrm{d}x^{2}$ and measure $\mathrm{d}x$ on
      $S^1_L$, is invariant under gauge transformations and under orientation-preserving
      diffeomorphisms, and is therefore a well-defined invariant relation on $\Sgeo$.
\end{enumerate}
\end{lemma}

\begin{proof}
Let $f$ be an orientation-preserving diffeomorphism of $S^1_L$ and let $A$ be a connection. Since
$f$ preserves the orientation, it fixes the fundamental generator $c=[\gamma_0]$ of
$H_1(S^1_L;\Int)$. Because $G=U(1)$ is abelian and every connection on $S^1_L$ is flat by
\eqref{eq:flat}, the holonomy of a connection around a loop depends on that loop only through its
class in $H_1(S^1_L;\Int)$; the pushforward $f_*A$ therefore has the same holonomy around $c$ as
$A$ does:
\begin{equation}
\label{eq:hol-diffeo}
\Hol\bigl(f_*A\bigr)=\Hol(A) .
\end{equation}
(Concretely, $\int_{S^1_L} f_*A=\int_{S^1_L}A$ by the change-of-variables formula for the
orientation-preserving $f$, the integrand being a $1$-form.)

\emph{(1).} By \eqref{eq:hol-diffeo} and the sufficiency direction of Theorem~\ref{thm:hol-class},
equality of holonomy already implies gauge equivalence; hence $f_*A$ is gauge equivalent to $A$ for
every orientation-preserving $f$ and every $A$. The equivalence relation generated by the gauge
group together with the orientation-preserving diffeomorphisms therefore coincides with gauge
equivalence alone, so the quotient \eqref{eq:circle-source-def} is unchanged by adjoining them, and
\eqref{eq:Sgeo-circle} holds.

\emph{(2).} Let $A$ and $A'$ be connections related either by a gauge transformation or by the
pushforward along an orientation-preserving diffeomorphism. In the second case part (1) supplies a
gauge transformation achieving the same result, so in either case there is $u\in\Gau(P)$ with
$A'=A^{u}$. The operator $\mathcal{D}$ is defined with respect to the fixed background metric and
measure of Convention~\ref{conv:circle-scope}, which are not transformed; the diffeomorphism
therefore plays no further role, and only the gauge covariance established in
Convention~\ref{conv:u1-norm} is needed. There we verified
$D_{A^{u}}\phi^{u}=u\,D_A\phi$ for $\phi^{u}=u\phi$; iterating,
\[
\mathcal{D}_{A^{u}}\bigl(u\,\phi\bigr)
=u\,\mathcal{D}_{A}\phi
\qquad\text{for all }\phi\in C^\infty(S^1_L,\Cx).
\]
Thus the multiplication operator $\mathcal{U}_u\phi:=u\phi$ intertwines $\mathcal{D}_A$ with
$\mathcal{D}_{A^{u}}$. Since $|u(x)|=1$ pointwise, $\mathcal{U}_u$ is an isometry of
$L^2(S^1_L,\mathrm{d}x)$ with inverse $\mathcal{U}_{u^{-1}}$, hence unitary; and it preserves
$C^\infty(S^1_L,\Cx)$, hence carries the domain of $\mathcal{D}_A$ onto that of
$\mathcal{D}_{A^{u}}$. Therefore $\mathcal{D}_{A}$ and $\mathcal{D}_{A'}=\mathcal{D}_{A^{u}}$ are
unitarily equivalent, and their spectral multisets coincide. Since $A$ and $A'$ were an arbitrary
pair related by a gauge transformation or by an allowed diffeomorphism, the spectral multiset is a
well-defined function on $\Sgeo$, which is the assertion.
\end{proof}

\begin{corollary}[A strictly coarser spectral completion]
\label{cor:spec}
Retain the source architecture of Corollary~\ref{cor:hol-completion}: the source is
$\Sgeo=\Mflat(S^1,U(1))$ and the declared Tier-1 readout remains the curvature $p_F$, which is
constant. Adjoin, as the invariant relation, the spectrum of the charged-scalar operator
\begin{equation}
\label{eq:Dop}
\mathcal{D}_A:=-\bigl(\partial_x-i\,a(x)\bigr)^{2}
\qquad\text{on}\qquad
C^\infty(S^1_L,\Cx)\subseteq L^2(S^1_L,\mathrm{d}x),
\end{equation}
computed with respect to the fixed standard metric $\mathrm{d}x^2$ and measure $\mathrm{d}x$ on
$S^1_L$, with the conventions of Convention~\ref{conv:u1-norm} and the standard self-adjoint
domain (the closure of $\mathcal{D}_A$ on $C^\infty(S^1_L,\Cx)$; $S^1_L$ is closed, so no boundary
conditions are needed and the operator is essentially self-adjoint there). The unitary charge-one
representation ensures that gauge transformations act unitarily on the charged-scalar Hilbert
space, so the spectral multiset is gauge invariant and hence a well-defined function on the declared
gauge quotient \eqref{eq:circle-source-def}. Its invariance under the auxiliary
orientation-preserving diffeomorphism action---which, as noted in Convention~\ref{conv:circle-scope},
is treated separately and does not consist of arrows of the fixed-metric groupoid, a general such
diffeomorphism not preserving $\mathrm{d}x$---is proved separately in
Lemma~\ref{lem:spec-invariance}(2). Set
$J_{\mathrm{spec}}([A]):=\operatorname{spec}\mathcal{D}_A$. Writing $\Hol(A)=e^{i\vartheta}$,
\begin{equation}
\label{eq:spec}
\operatorname{spec}\mathcal{D}_A
=\left\{
\left(\frac{2\pi k-\vartheta}{L}\right)^{\!2}
:\; k\in\Int
\right\}
\qquad\text{(as a multiset)} ,
\end{equation}
and the resulting realized completion is
\begin{equation}
\label{eq:Espec}
E_{\mathrm{spec}}
\;\cong\;
U(1)\big/\bigl(z\sim z^{-1}\bigr)
\;\cong\;
[-1,1],
\end{equation}
which is \emph{strictly coarser} than $E_{\mathrm{hol}}$: the completion map
$\Sgeo\to E_{\mathrm{spec}}$ is surjective but not injective, its nontrivial fibres being exactly
the pairs $\{[A_\vartheta],[A_{-\vartheta}]\}$ with $\vartheta\not\equiv 0,\pi\pmod{2\pi}$.
\end{corollary}

\begin{proof}
\emph{The spectrum.} By Theorem~\ref{thm:hol-class} every class has the representative
$A_\vartheta=(\vartheta/L)\,\mathrm{d}x$ of \eqref{eq:Atheta}, and by
Example~\ref{ex:spectral}(ii) the spectrum is unchanged along the gauge orbit, so we may compute
with $A_\vartheta$. With $a\equiv\vartheta/L$ constant, $\mathcal{D}_{A_\vartheta}$ has the
complete orthonormal system of eigenfunctions $\psi_k(x)=L^{-1/2}e^{2\pi i k x/L}$, $k\in\Int$
(complete in $L^2(S^1_L,\mathrm{d}x)$ by Fourier theory). By the convention \eqref{eq:u1-covder},
\[
\Bigl(\partial_x-i\tfrac{\vartheta}{L}\Bigr)\psi_k
=\Bigl(\tfrac{2\pi i k}{L}-\tfrac{i\vartheta}{L}\Bigr)\psi_k
=i\,\frac{2\pi k-\vartheta}{L}\,\psi_k ,
\]
and hence, squaring and inserting the overall minus sign of \eqref{eq:Dop},
\[
\mathcal{D}_{A_\vartheta}\psi_k
=-\left(i\,\frac{2\pi k-\vartheta}{L}\right)^{\!2}\psi_k
=\left(\frac{2\pi k-\vartheta}{L}\right)^{\!2}\psi_k .
\]
Since the $\psi_k$ are complete, the spectrum is purely discrete and is exactly the multiset
\eqref{eq:spec}. Note that $\vartheta$ is defined only modulo $2\pi$, and \eqref{eq:spec} is
invariant under $\vartheta\mapsto\vartheta+2\pi$ via the relabelling $k\mapsto k+1$; the multiset
is therefore a well-defined function of $[A]$, as it must be.

\emph{``If'': inversion-related holonomies are spectrally equal.} Replacing $\vartheta$ by
$-\vartheta$ in \eqref{eq:spec} and relabelling $k\mapsto -k$ carries the multiset to itself, since
$(2\pi(-k)-(-\vartheta))^2=(-(2\pi k-\vartheta))^2=(2\pi k-\vartheta)^2$. Hence
$[A_\vartheta]$ and $[A_{-\vartheta}]$---which by Theorem~\ref{thm:hol-class} are distinct classes
whenever $e^{i\vartheta}\neq e^{-i\vartheta}$, i.e.\ whenever
$\vartheta\not\equiv0,\pi\pmod{2\pi}$---have the same spectrum.

\emph{``Only if'': equal spectra force inversion-related holonomies.} We prove this directly rather
than by citation. Suppose $\operatorname{spec}\mathcal{D}_{A_\vartheta}
=\operatorname{spec}\mathcal{D}_{A_{\vartheta'}}$ as multisets. Normalize
$\vartheta,\vartheta'\in[0,2\pi)$, which is permissible by the periodicity just noted. By
\eqref{eq:spec} every eigenvalue of $\mathcal{D}_{A_\vartheta}$ is a nonnegative real number, so we
may form the set of both real square roots of the eigenvalues, rescaled by $L/2\pi$:
\begin{align*}
\Sigma(\vartheta)
&:=\Bigl\{\pm\tfrac{L}{2\pi}\sqrt{\lambda}\;:\;
\lambda\in\operatorname{spec}\mathcal{D}_{A_\vartheta}\Bigr\}
\\
&\phantom{:}=
\Bigl\{k-\tfrac{\vartheta}{2\pi}\;:\;k\in\Int\Bigr\}
\;\cup\;
\Bigl\{-k+\tfrac{\vartheta}{2\pi}\;:\;k\in\Int\Bigr\}
\\
&\phantom{:}=
\Bigl(-\tfrac{\vartheta}{2\pi}+\Int\Bigr)
\;\cup\;
\Bigl(\tfrac{\vartheta}{2\pi}+\Int\Bigr),
\end{align*}
which is determined by the spectral multiset alone. Thus equality of spectra gives
\[
\Bigl(\tfrac{\vartheta}{2\pi}+\Int\Bigr)\cup\Bigl(-\tfrac{\vartheta}{2\pi}+\Int\Bigr)
\;=\;
\Bigl(\tfrac{\vartheta'}{2\pi}+\Int\Bigr)\cup\Bigl(-\tfrac{\vartheta'}{2\pi}+\Int\Bigr)
\]
as subsets of $\Real$.
Intersecting both sides with the interval $[0,1)$ and comparing, the unordered pair of residues
$\{\tfrac{\vartheta}{2\pi}\bmod 1,\;-\tfrac{\vartheta}{2\pi}\bmod 1\}$ equals the pair
$\{\tfrac{\vartheta'}{2\pi}\bmod 1,\;-\tfrac{\vartheta'}{2\pi}\bmod 1\}$. Hence
$\vartheta'\equiv\vartheta$ or $\vartheta'\equiv-\vartheta \pmod{2\pi}$, that is,
$\Hol(A')=\Hol(A)^{\pm1}$, as claimed.

\emph{The completion.} By the two directions just proved, the relation $J_{\mathrm{spec}}$ has
exactly the level sets of the map $U(1)\to U(1)/(z\sim z^{-1})$. Since the readout $p_F$ is constant
and contributes nothing, Definition~\ref{def:egeo} gives
$E_{\mathrm{spec}}=\Imm(p_F,J_{\mathrm{spec}})\cong U(1)/(z\sim z^{-1})$. The map
$z=e^{i\vartheta}\mapsto\tfrac12(z+z^{-1})=\cos\vartheta$ is a continuous bijection from this
quotient onto $[-1,1]$ (it is well defined and injective on the quotient, and surjective onto
$[-1,1]$), and it is a homeomorphism since the quotient is compact and $[-1,1]$ is Hausdorff. This
gives \eqref{eq:Espec}. Non-injectivity of the completion map and the description of its fibres now
follow from Proposition~\ref{prop:orbit-reconstruction} applied to the coarser relation family
$\{p_F\}\cup\{J_{\mathrm{spec}}\}$.
\end{proof}

\begin{remark}[The two corollaries together]
\label{rem:two-corollaries}
Corollaries~\ref{cor:hol-completion} and \ref{cor:spec} are two halves of one lesson, and both are
consequences of Theorem~\ref{thm:restricted-realization} and
Proposition~\ref{prop:orbit-reconstruction} applied to the \emph{same} source $\Sgeo$ with the
\emph{same} constant readout $p_F$; only the declared invariant relation differs. The realized
completion is relative to the declared relation family (Remark~\ref{rem:relative}). Declare holonomy
and the source is reconstructed exactly; declare instead the charged-scalar spectrum and a residual
$z\sim z^{-1}$ degeneracy survives the completion---the shadow of the complex conjugation
$\phi\mapsto\bar\phi$, equivalently of the orientation reversal excluded in
Definition~\ref{def:circle-model} (Remark~\ref{rem:u1-hyps}). This is the Paper~2 phenomenon
\cite{ST2}, realized geometrically: whether the lost information obstructs the nominated target
depends on the target.
\end{remark}

% ======================================================================
\section{Variational realization of completed relations}
\label{sec:variational}

We now embed the completed relation data into dynamics. The claim to be established is narrow and
precise: an invariant relation-dependent action yields, by functional variation, correctly
normalized stress-energy, gauge-current, and matter responses, and nothing is postulated by naming
it.

\subsection{Variational setting}

\begin{assumption}[Well-posed variational problem]
\label{ass:variational}
The global configuration satisfies the descent conditions of Theorem~\ref{thm:descent} and lies in
$\Cgeo$. Furthermore, \emph{one} of the following holds, and which one is part of the declaration of
the model:
\begin{enumerate}[label=(V\arabic*),leftmargin=2.8em]
\item\label{V1} all variations $\delta g$, $\delta A$, $\delta\Phi$ have compact support in the
      interior of $M$; or
\item\label{V2} boundary values of the fields, and for gravity the induced metric on $\partial M$,
      are held fixed, so that boundary contributions vanish identically; or
\item\label{V3} the gravitational and gauge actions are supplemented by the boundary terms required
      for a well-posed variational problem. For gravity this is the Gibbons--Hawking--York
      term
      \[
      \frac{1}{8\pi G_N}\oint_{\partial M} K\,\sqrt{|h|}\;\mathrm{d}^3y
      \]
      \cite{York1972,GibbonsHawking1977}, with the appropriate subtraction or counterterms in
      the asymptotic case; for the gauge sector, whatever boundary term the declared $\mathcal{B}$
      requires.
\end{enumerate}
All functional derivatives below are asserted to exist on the declared domain; where
$S_{\mathrm{rel}}$ is nonlocal, this existence is a hypothesis on the model, not a theorem.
\end{assumption}

\begin{definition}[Relation-dependent total action]
\label{def:action}
Let $r=(r_\alpha)_{\alpha\in A}$ denote the completed invariant relation coordinates, i.e.\ the
components of $\egeo$. Consider
\begin{equation}
\label{eq:Stot}
S_{\mathrm{tot}}[g,A,\Phi;r]
=
S_{\mathrm{EH}}[g]
+S_{\mathrm{YM}}[g,A]
+S_{\mathrm{m}}[g,A,\Phi]
+S_{\mathrm{rel}}[g,A,\Phi;r],
\end{equation}
with
\begin{align}
\label{eq:SEH}
S_{\mathrm{EH}}[g]
&=\frac{1}{16\pi G_N}\int_M \bigl(R-2\Lambda\bigr)\,\dvol
\;\;(+\text{ boundary term under \ref{V3}}),
\\[0.3em]
\label{eq:SYM}
S_{\mathrm{YM}}[g,A]
&=-\frac{1}{4g_{\mathrm{YM}}^{2}}\int_M
\bigl\langle F_{\mu\nu},F^{\mu\nu}\bigr\rangle\,\dvol ,
\end{align}
and $S_{\mathrm{m}}$ a matter action for $\Phi$ \cite{YangMills1954}. The functional
$S_{\mathrm{rel}}$ carries the model-specific dependence on $r$. It may be local,
boundary-supported, or nonlocal, provided that
\begin{enumerate}[label=(R\arabic*),leftmargin=2.8em]
\item\label{R1} it depends on the relation coordinates through a finite or countable dynamical
      subfamily $A_{\mathrm{dyn}}\subseteq A$, and its Fr\'echet derivatives with respect to
      $g^{\mu\nu}$, $A_\mu$, and the independent matter components exist on the stated domain and
      are represented by distributional kernels, as set out in Section~\ref{sec:chain-rule}; and
\item\label{R2} it is invariant under the allowed bundle automorphisms and allowed
      diffeomorphisms---that is, $S_{\mathrm{rel}}$ is constant on the arrows of $\Ggeo$, exactly as
      the $J_\alpha$ are (Definition~\ref{def:invariant-relation}).
\end{enumerate}
\end{definition}

Condition \ref{R2} is what makes $S_{\mathrm{rel}}$ a functional of the \emph{physical orbit}, and it
is what supplies the Ward identities of Section~\ref{sec:noether}. An $S_{\mathrm{rel}}$ that fails to be $\Ggeo$-invariant does not define a geometric realization of
a completed relation.

\subsection{The functional chain rule for relation functionals}
\label{sec:chain-rule}

\begin{convention}[Two regimes for $r$, and the correct chain rule]
\label{conv:two-regimes}
By Convention~\ref{conv:relation-status}, exactly one of the following applies to each component
$r_\alpha$, and the model declares which.

\medskip
\noindent\textbf{(a) $r_\alpha$ is a functional of the fields.}
Paper~3 permits an arbitrary index \emph{set} $A$, and nothing in
Sections~\ref{sec:realization}--\ref{sec:u1} requires more. A differentiable action, however, needs
additional analytic structure, and we impose it only where it is used. We therefore require that
$S_{\mathrm{rel}}$ depend on the relation coordinates through a distinguished \emph{dynamical
subfamily}
\begin{equation}
\label{eq:Adyn}
A_{\mathrm{dyn}}\subseteq A,
\qquad
A_{\mathrm{dyn}} \text{ finite or countable},
\end{equation}
so that $S_{\mathrm{rel}}=S_{\mathrm{rel}}[g,A,\Phi;(r_\alpha)_{\alpha\in A_{\mathrm{dyn}}}]$ and
$\partial S_{\mathrm{rel}}/\partial r_\alpha=0$ for $\alpha\notin A_{\mathrm{dyn}}$. When
$A_{\mathrm{dyn}}$ is countably infinite we require in addition that the derivative series in
\eqref{eq:chain-rule} below converge \emph{unconditionally} in the distributional-dual topology of
the relevant space of kernels; when $A_{\mathrm{dyn}}$ is finite this is vacuous. The remaining
relation coordinates, indexed by $A\setminus A_{\mathrm{dyn}}$, continue to serve as completion
coordinates in the sense of Sections~\ref{sec:realization}--\ref{sec:u1}; they simply do not enter
the action.

Now suppose $r_\alpha=r_\alpha[g,A,\Phi]$, $\alpha\in A_{\mathrm{dyn}}$, takes values in a Banach
(or finite-dimensional) space $V_\alpha$ and is Fr\'echet differentiable on the declared domain.
Then $r_\alpha$ is in general \emph{nonlocal}: a Wilson observable integrates around a loop, an
integral invariant integrates over $M$, a spectral quantity depends on the fields through a global
eigenvalue problem. Its variation is accordingly given not by a pointwise partial derivative but by
the pairing of distributional kernels with the variations:
\begin{equation}
\label{eq:drho}
\delta r_\alpha
=\int_M\!\Bigl[
\frac{\delta r_\alpha}{\delta g^{\mu\nu}(y)}\,\delta g^{\mu\nu}(y)
+\Bigl\langle \frac{\delta r_\alpha}{\delta A_\mu(y)},\,\delta A_\mu(y)\Bigr\rangle
+\Bigl\langle\!\Bigl\langle \frac{\delta r_\alpha}{\delta \Phi(y)},\,\delta\Phi(y)
\Bigr\rangle\!\Bigr\rangle
\Bigr]\sqrt{-g(y)}\;\mathrm{d}^4y ,
\end{equation}
where $\delta r_\alpha/\delta g^{\mu\nu}(y)$ etc.\ denote the $V_\alpha$-valued distributional
kernels representing the derivative, $\langle\cdot,\cdot\rangle$ is the $\Ad$-invariant pairing on
$\gfrak$, and $\langle\!\langle\cdot,\cdot\rangle\!\rangle$ is the duality pairing between
$W^{*}$ and $W$ discussed in Section~\ref{sec:matter-pairing}. (For a Wilson loop $r_\alpha$ the
kernel $\delta r_\alpha/\delta A_\mu(y)$ is supported on the loop and is a distribution, not a
function; \eqref{eq:drho} is exactly the statement that must be used, and a pointwise
$\partial r_\alpha/\partial A_\mu(y)$ does not exist.)

Then, writing $\partial S_{\mathrm{rel}}/\partial r_\alpha\in V_\alpha^{*}$ for the derivative of
$S_{\mathrm{rel}}$ with respect to its $V_\alpha$-argument (a bounded linear functional, and an
ordinary partial derivative only when $V_\alpha$ is finite-dimensional), the chain rule reads
\begin{equation}
\label{eq:chain-rule}
\frac{\delta S_{\mathrm{rel}}}{\delta g^{\mu\nu}(y)}
=\left.\frac{\delta S_{\mathrm{rel}}}{\delta g^{\mu\nu}(y)}\right|_{r\;\mathrm{fixed}}
+\sum_{\alpha\in A_{\mathrm{dyn}}}
\Bigl\langle \frac{\partial S_{\mathrm{rel}}}{\partial r_\alpha},\;
\frac{\delta r_\alpha}{\delta g^{\mu\nu}(y)}\Bigr\rangle_{V_\alpha^*\!,V_\alpha} ,
\end{equation}
and identically with $g^{\mu\nu}(y)$ replaced by $A_\mu(y)$ or by the matter components. The sum is
over the dynamical subfamily \eqref{eq:Adyn}, and converges by the unconditional-convergence
hypothesis imposed there. Only in the special case where $V_\alpha=\Real$ and $r_\alpha$ is a
\emph{local} scalar field does \eqref{eq:chain-rule} reduce to an ordinary pointwise partial
derivative; in general it does not.

\medskip
\noindent\textbf{(b) $r_\alpha$ labels a fixed topological or boundary sector.}
Then $\delta r_\alpha=0$ under all variations within the sector---this is what ``fixed sector''
means---and $r_\alpha$ enters the field equations only as a parameter. The second term of
\eqref{eq:chain-rule} is absent.

\medskip
In neither case is $r_\alpha$ varied as an independent dynamical field. A model wishing to promote
$r_\alpha$ to a field must adjoin it to the configuration tuple, enlarge $\Cgeo$, and supply its own
kinetic term; that is a different model, and its Euler--Lagrange equation is an additional equation
not derived here.
\end{convention}

\subsection{Matter pairings and independent components}
\label{sec:matter-pairing}

\begin{convention}[Matter variation and the dual pairing]
\label{conv:matter}
The representation $\rho\colon G\to\mathrm{GL}(W)$ is not assumed unitary
(Convention~\ref{conv:basic}), and $W$ may be complex. Accordingly the matter Euler operator is
valued in the dual, and the variational pairing must be specified. We adopt the following, and use
it consistently.

\emph{(i) Real independent components.} Suppose $\Phi$ denotes a collection of \emph{real}
independent field components. This is always achievable: decompose a complex field into its real
and imaginary parts, or a section of a complex bundle into its underlying real bundle. Then
$\delta S/\delta\Phi$ is valued in the real dual bundle $P\times_{\rho^{*}}W^{*}$, and the
pairing
\[
\langle\!\langle\cdot,\cdot\rangle\!\rangle\colon W^{*}\times W\longrightarrow\Real
\]
is the canonical duality pairing. This pairing is $G$-invariant for the contragredient action
$\rho^{*}(h)=\rho(h^{-1})^{\mathsf{T}}$, and requires no metric on $W$. There is then a single
Euler equation $\mathcal{E}_\Phi=0$.

\emph{(ii) Complex fields treated holomorphically.} If instead $\Phi$ is treated as a complex field
with $\Phi$ and $\bar\Phi$ varied independently---the usual convention for a charged scalar or a
Dirac spinor---then there are \emph{two} Euler equations,
\begin{equation}
\label{eq:conj-euler}
\mathcal{E}_{\Phi}:=\frac{1}{\sqrt{-g}}\frac{\delta S}{\delta \bar\Phi}=0
\qquad\text{and}\qquad
\overline{\mathcal{E}_{\Phi}}:=\frac{1}{\sqrt{-g}}\frac{\delta S}{\delta \Phi}=0 ,
\end{equation}
which are conjugate to one another when the action is real, and the pairing used in the variation is
the sesquilinear one furnished by a $\rho$-invariant Hermitian form on $W$---which exists precisely
when $\rho$ is unitary. If $\rho$ is not unitary, no such Hermitian form is available and one must
revert to convention (i).

For definiteness, and to avoid carrying two conjugate equations through every formula, \emph{in the
statements below $\Phi$ denotes real independent field components in the sense of (i)}, and
$\mathcal{E}_\Phi$ is the corresponding single dual-valued Euler operator. Every result restates
verbatim in convention (ii) upon replacing $\mathcal{E}_{\Phi}=0$ by the conjugate pair
\eqref{eq:conj-euler} and $\langle\!\langle\cdot,\cdot\rangle\!\rangle$ by the invariant Hermitian
pairing.
\end{convention}

\subsection{Response tensors, currents, and Euler operators}

\begin{definition}[Responses]
\label{def:responses}
For each nongravitational sector $X\in\{\mathrm{YM},\mathrm{m},\mathrm{rel}\}$ define the
stress-energy tensor
\begin{equation}
\label{eq:Tdef}
T^{X}_{\mu\nu}
\;:=\;
-\frac{2}{\sqrt{-g}}\,\frac{\delta S_X}{\delta g^{\mu\nu}} ,
\end{equation}
the gauge current
\begin{equation}
\label{eq:Jdef}
J^{\mu}_{\mathrm{m}}
\;:=\;
-\frac{1}{\sqrt{-g}}\,\frac{\delta S_{\mathrm{m}}}{\delta A_\mu},
\qquad
J^{\mu}_{\mathrm{rel}}
\;:=\;
-\frac{1}{\sqrt{-g}}\,\frac{\delta S_{\mathrm{rel}}}{\delta A_\mu},
\end{equation}
each valued in $\gfrak$ (with $\gfrak$-indices contracted using $\langle\cdot,\cdot\rangle$), and the
matter Euler operators
\begin{equation}
\label{eq:Edef}
\mathcal{E}^{X}_{\Phi}
\;:=\;
\frac{1}{\sqrt{-g}}\,\frac{\delta S_X}{\delta\Phi} ,
\end{equation}
valued in the dual bundle and paired with $\delta\Phi$ by
$\langle\!\langle\cdot,\cdot\rangle\!\rangle$ as in Convention~\ref{conv:matter}(i). For
$X=\mathrm{rel}$ the functional derivatives in \eqref{eq:Tdef}--\eqref{eq:Edef} are those computed by
the chain rule \eqref{eq:chain-rule}, and thus already include the relation-mediated contributions.
\end{definition}

\begin{lemma}[Yang--Mills stress tensor]
\label{lem:TYM}
With the normalization \eqref{eq:SYM},
\begin{equation}
\label{eq:TYM}
T^{\mathrm{YM}}_{\mu\nu}
=\frac{1}{g_{\mathrm{YM}}^{2}}
\left(
\bigl\langle F_{\mu\rho},F_{\nu}{}^{\rho}\bigr\rangle
-\tfrac14\,g_{\mu\nu}\,\bigl\langle F_{\rho\sigma},F^{\rho\sigma}\bigr\rangle
\right),
\end{equation}
which is symmetric, gauge invariant, and (in four dimensions) trace free.
\end{lemma}

\begin{proof}
Write $S_{\mathrm{YM}}=-\tfrac{1}{4g_{\mathrm{YM}}^2}\int_M
\langle F_{\mu\rho},F_{\nu\sigma}\rangle g^{\mu\nu}g^{\rho\sigma}\sqrt{-g}\,\mathrm{d}^4x$; the
components $F_{\mu\nu}=\partial_\mu A_\nu-\partial_\nu A_\mu+[A_\mu,A_\nu]$ are independent of $g$.
Using $\delta\sqrt{-g}=-\tfrac12\sqrt{-g}\,g_{\mu\nu}\,\delta g^{\mu\nu}$ and varying the two inverse
metrics,
\[
\delta S_{\mathrm{YM}}
=-\frac{1}{4g_{\mathrm{YM}}^2}\int_M
\Bigl(
2\langle F_{\mu\rho},F_{\nu}{}^{\rho}\rangle
-\tfrac12 g_{\mu\nu}\langle F_{\rho\sigma},F^{\rho\sigma}\rangle
\Bigr)\delta g^{\mu\nu}\sqrt{-g}\,\mathrm{d}^4x ,
\]
and \eqref{eq:Tdef} gives \eqref{eq:TYM}. Gauge invariance follows from $\Ad$-invariance of
$\langle\cdot,\cdot\rangle$ and \eqref{eq:gauge-transf}; tracelessness from
$g^{\mu\nu}\langle F_{\mu\rho},F_\nu{}^\rho\rangle=\langle F_{\rho\sigma},F^{\rho\sigma}\rangle$ and
$g^{\mu\nu}g_{\mu\nu}=4$.
\end{proof}

\begin{lemma}[Yang--Mills variation with respect to the connection]
\label{lem:YMvar}
Under Assumption~\ref{ass:variational},
\begin{equation}
\label{eq:dSYM}
\delta S_{\mathrm{YM}}
=-\frac{1}{g_{\mathrm{YM}}^{2}}\int_M
\bigl\langle F^{\mu\nu},\,D_\mu\,\delta A_\nu\bigr\rangle\,\dvol
=\;\frac{1}{g_{\mathrm{YM}}^{2}}\int_M
\bigl\langle D_\mu F^{\mu\nu},\,\delta A_\nu\bigr\rangle\,\dvol ,
\end{equation}
up to the treated boundary term.
\end{lemma}

\begin{proof}
From $F_{\mu\nu}=\partial_\mu A_\nu-\partial_\nu A_\mu+[A_\mu,A_\nu]$ one has
\begin{equation}
\label{eq:dF}
\delta F_{\mu\nu}=D_\mu\,\delta A_\nu - D_\nu\,\delta A_\mu .
\end{equation}
Hence, from \eqref{eq:SYM},
\[
\delta S_{\mathrm{YM}}
=-\frac{1}{4g_{\mathrm{YM}}^{2}}\int_M
2\,\bigl\langle F^{\mu\nu},\,\delta F_{\mu\nu}\bigr\rangle\,\dvol
=-\frac{1}{2g_{\mathrm{YM}}^{2}}\int_M
\bigl\langle F^{\mu\nu},\,D_\mu\delta A_\nu - D_\nu\delta A_\mu\bigr\rangle\,\dvol .
\]
The two terms in the bracket are \emph{equal} after relabelling, by the antisymmetry of $F^{\mu\nu}$:
swapping the dummy indices $\mu\leftrightarrow\nu$ in the second term gives
$-\langle F^{\nu\mu},D_\mu\delta A_\nu\rangle=+\langle F^{\mu\nu},D_\mu\delta A_\nu\rangle$.
Therefore the bracket contributes twice the first term, the factor $\tfrac12$ cancels, and
\[
\delta S_{\mathrm{YM}}
=-\frac{1}{g_{\mathrm{YM}}^{2}}\int_M
\bigl\langle F^{\mu\nu},\,D_\mu\delta A_\nu\bigr\rangle\,\dvol ,
\]
which is the first equality in \eqref{eq:dSYM}. Covariant integration by parts---legitimate
because
$\langle\cdot,\cdot\rangle$ is $\Ad$-invariant, hence $D$-compatible, and because the boundary term
vanishes by Assumption~\ref{ass:variational}---moves $D_\mu$ onto $F^{\mu\nu}$ and flips the sign,
giving the second equality.
\end{proof}

\subsection{The coupled field equations}

\begin{theorem}[Coupled field equations]
\label{thm:field-eqs}
Under Assumption~\ref{ass:variational} and Conventions~\ref{conv:two-regimes} and
\ref{conv:matter}(i), stationarity of $S_{\mathrm{tot}}$ under independent variations of
$g^{\mu\nu}$, $A_\mu$, and $\Phi$ is equivalent to the system
\begin{align}
\label{eq:einstein}
G_{\mu\nu}+\Lambda g_{\mu\nu}
&=8\pi G_N\Bigl(
T^{\mathrm{YM}}_{\mu\nu}+T^{\mathrm{m}}_{\mu\nu}+T^{\mathrm{rel}}_{\mu\nu}
\Bigr),
\\[0.35em]
\label{eq:ym}
\frac{1}{g_{\mathrm{YM}}^{2}}\,D_\mu F^{\mu\nu}
&=J^{\nu}_{\mathrm{m}}+J^{\nu}_{\mathrm{rel}},
\\[0.35em]
\label{eq:matter}
\mathcal{E}^{\mathrm{m}}_{\Phi}+\mathcal{E}^{\mathrm{rel}}_{\Phi}
&=0,
\end{align}
with representation indices and the pairings $\langle\cdot,\cdot\rangle$ and
$\langle\!\langle\cdot,\cdot\rangle\!\rangle$ restored as required by $G$ and $\rho$.
\end{theorem}

\begin{proof}
\emph{Metric variation.}
The standard variation of \eqref{eq:SEH}, after the boundary treatment stipulated in
Assumption~\ref{ass:variational} (which is exactly what disposes of the total-derivative term
$\int_M g^{\mu\nu}\delta R_{\mu\nu}\sqrt{-g}$; under \ref{V3} it is the Gibbons--Hawking--York term
that cancels it \cite{York1972,GibbonsHawking1977}), gives
\begin{equation}
\label{eq:dSEH}
\delta S_{\mathrm{EH}}
=\frac{1}{16\pi G_N}\int_M
\bigl(G_{\mu\nu}+\Lambda g_{\mu\nu}\bigr)\,\delta g^{\mu\nu}\,\dvol .
\end{equation}
By Definition~\ref{def:responses},
\begin{equation}
\label{eq:dSmatter}
\delta\bigl(S_{\mathrm{YM}}+S_{\mathrm{m}}+S_{\mathrm{rel}}\bigr)\Big|_{\delta g}
=-\frac12\int_M
\bigl(
T^{\mathrm{YM}}_{\mu\nu}+T^{\mathrm{m}}_{\mu\nu}+T^{\mathrm{rel}}_{\mu\nu}
\bigr)\,\delta g^{\mu\nu}\,\dvol ,
\end{equation}
where $\delta S_{\mathrm{rel}}$ is computed by the chain rule \eqref{eq:chain-rule}: in regime (a)
of Convention~\ref{conv:two-regimes} it includes the relation-mediated kernel contributions, and in
regime (b) those are absent. In either case \eqref{eq:dSmatter} \emph{defines} its left side via
\eqref{eq:Tdef} and adds no assumption. Setting $\delta S_{\mathrm{tot}}=0$ for arbitrary symmetric
$\delta g^{\mu\nu}$ (compactly supported, or boundary-vanishing) and applying the fundamental lemma
of the calculus of variations gives
\[
\frac{1}{16\pi G_N}\bigl(G_{\mu\nu}+\Lambda g_{\mu\nu}\bigr)
-\frac12\bigl(T^{\mathrm{YM}}_{\mu\nu}+T^{\mathrm{m}}_{\mu\nu}+T^{\mathrm{rel}}_{\mu\nu}\bigr)=0 ,
\]
which is \eqref{eq:einstein} with the displayed $8\pi G_N$ normalization.

\emph{Gauge variation.}
By Lemma~\ref{lem:YMvar},
\[
\delta S_{\mathrm{YM}}\big|_{\delta A}
=\frac{1}{g_{\mathrm{YM}}^{2}}\int_M
\bigl\langle D_\mu F^{\mu\nu},\,\delta A_\nu\bigr\rangle\,\dvol ,
\]
while by the definitions \eqref{eq:Jdef},
\[
\delta\bigl(S_{\mathrm{m}}+S_{\mathrm{rel}}\bigr)\big|_{\delta A}
=-\int_M\bigl\langle J^{\nu}_{\mathrm{m}}+J^{\nu}_{\mathrm{rel}},\,
\delta A_\nu\bigr\rangle\,\dvol .
\]
Setting the sum of the two to zero for arbitrary $\delta A_\nu$, and using the nondegeneracy of
$\langle\cdot,\cdot\rangle$, gives \eqref{eq:ym}.

\emph{Matter variation.}
$S_{\mathrm{EH}}$ and $S_{\mathrm{YM}}$ are independent of $\Phi$, so by \eqref{eq:Edef}
$\delta_\Phi S_{\mathrm{tot}}
=\int_M \langle\!\langle\,\mathcal{E}^{\mathrm{m}}_\Phi+\mathcal{E}^{\mathrm{rel}}_\Phi,\;
\delta\Phi\,\rangle\!\rangle\,\dvol$; arbitrariness of the real independent components $\delta\Phi$
and nondegeneracy of the duality pairing give \eqref{eq:matter}. Under
Convention~\ref{conv:matter}(ii) this is replaced by the conjugate pair \eqref{eq:conj-euler}.

The three variations are independent, so stationarity of $S_{\mathrm{tot}}$ is equivalent to the
simultaneous validity of \eqref{eq:einstein}--\eqref{eq:matter}.
\end{proof}

\begin{remark}[What has and has not been derived]
\label{rem:not-postulated}
The tensor $T^{\mathrm{rel}}_{\mu\nu}$ is not an independent postulate. By \eqref{eq:Tdef} and the
chain rule \eqref{eq:chain-rule} it is the metric functional derivative of a specified
$\Ggeo$-invariant action satisfying \ref{R1}--\ref{R2}. Different source completions, and
different couplings of the same completion, produce different $T^{\mathrm{rel}}_{\mu\nu}$,
$J^{\mu}_{\mathrm{rel}}$, $\mathcal{E}^{\mathrm{rel}}_{\Phi}$. Theorem~\ref{thm:field-eqs} asserts
only the \emph{structure} of the response---that a completed relation enters Tier-1 dynamics through
functional variation of an invariant action, with the normalizations displayed---and asserts nothing
about the magnitude, sign, or empirical existence of any such term. No new interaction is proposed,
and no empirical claim is made.
\end{remark}

% ======================================================================
\section{Ward--Noether identities and covariant consistency}
\label{sec:noether}

Gauge and diffeomorphism invariance impose Noether identities \cite{Noether1918}, and those
identities constrain every effective representation of a completed source relation. Two logically
distinct routes to covariant conservation must be kept apart, and we present them separately.

\begin{convention}[Infinitesimal allowed arrows and the induced field variations]
\label{conv:lie}
A base vector field $\xi$ on $M$ does not by itself define a Lie derivative of a principal
connection or of a section of an associated bundle: differentiating such objects requires a lift of
the flow of $\xi$ to the total space $P$. We therefore generate the variations from the groupoid
itself.

Let $(\hat F_t,f_t)_{t\in(-\epsilon,\epsilon)}$ be a smooth one-parameter family of allowed arrows
of $\Ggeo$ with $(\hat F_0,f_0)=(\id_P,\id_M)$, and let
\[
\xi:=\frac{\mathrm{d}}{\mathrm{d}t}\Big|_{t=0} f_t
\]
be the generating base vector field, taken compactly supported in the interior of $M$. Write
$\hat\xi$ for the induced $G$-invariant vector field on $P$ generating $\hat F_t$, which projects to
$\xi$. The induced infinitesimal variations of the fields are defined by
\begin{equation}
\label{eq:delta-xi}
\begin{aligned}
\delta_\xi g^{\mu\nu}
&:=-\frac{\mathrm{d}}{\mathrm{d}t}\Big|_{t=0}\bigl(f_t^{*}g^{\mu\nu}\bigr)
=-\Lie_\xi g^{\mu\nu},
\\[0.35em]
\delta_{\hat\xi} A
&:=-\frac{\mathrm{d}}{\mathrm{d}t}\Big|_{t=0}\bigl(\hat F_t^{*}A\bigr),
\\[0.35em]
\delta_{\hat\xi}\Phi
&:=-\frac{\mathrm{d}}{\mathrm{d}t}\Big|_{t=0}
\bigl(\hat F_{t,\rho}^{-1}\circ\Phi\circ f_t\bigr).
\end{aligned}
\end{equation}
The first is intrinsic to the base, and we record the inverse-metric identity we shall use:
differentiating $g^{\mu\rho}g_{\rho\nu}=\delta^{\mu}_{\nu}$ gives
\begin{equation}
\label{eq:lie-inverse}
\Lie_\xi g^{\mu\nu}=-\bigl(\nabla^{\mu}\xi^{\nu}+\nabla^{\nu}\xi^{\mu}\bigr),
\end{equation}
the sign being opposite to that for the covariant metric,
$\Lie_\xi g_{\mu\nu}=\nabla_\mu\xi_\nu+\nabla_\nu\xi_\mu$. Hence, with the convention
\eqref{eq:delta-xi},
\begin{equation}
\label{eq:delta-xi-metric}
\delta_\xi g^{\mu\nu}=\nabla^{\mu}\xi^{\nu}+\nabla^{\nu}\xi^{\mu} .
\end{equation}

The bundle variations require a further remark. The lift $\hat\xi$ of $\xi$ is not unique: any two
allowed lifts differ by a vertical, i.e.\ gauge, direction. Concretely, in a local trivialization
over $U\subseteq M$, in which $A$ is represented by $A_\mu$ and $\Phi$ by a $W$-valued map, a choice
of lift yields
\begin{equation}
\label{eq:delta-xi-local}
\delta_{\hat\xi}A_\mu = -\bigl(\Lie_\xi A\bigr)_\mu + D_\mu\varepsilon,
\qquad
\delta_{\hat\xi}\Phi = -\Lie_\xi\Phi - \rho_*(\varepsilon)\,\Phi ,
\end{equation}
where $\Lie_\xi$ denotes the coordinate Lie derivative of the local representatives and
$\varepsilon\in\Gamma(\mathrm{ad}P|_U)$ is the vertical gauge parameter recording the choice of
lift; changing the lift changes $\varepsilon$. Thus the induced connection variation agrees with the
naive coordinate Lie derivative only up to a vertical term $D_\mu\varepsilon$, with the
corresponding representation action $-\rho_*(\varepsilon)\Phi$ on the matter field.

This ambiguity is harmless for the Ward identity of Theorem~\ref{thm:cons}, and we note why once
here. The extra terms in \eqref{eq:delta-xi-local} are precisely an infinitesimal gauge
transformation, of the form \eqref{eq:inf-gauge}. In the Ward computation they are contracted with
the Euler coefficients of $A$ and $\Phi$, which vanish identically on the gauge and matter shell on
which the theorem is evaluated. The vertical terms therefore multiply zero and drop out, and the
conclusion is independent of the choice of lift.
\end{convention}

\begin{definition}[Total stress-energy]
\label{def:Ttot}
\begin{equation}
\label{eq:Ttot}
T^{\mathrm{tot}}_{\mu\nu}
:=T^{\mathrm{YM}}_{\mu\nu}+T^{\mathrm{m}}_{\mu\nu}+T^{\mathrm{rel}}_{\mu\nu} .
\end{equation}
\end{definition}

\subsection{Route (a): on a full solution, via the contracted Bianchi identity}

\begin{proposition}[Conservation on a full solution]
\label{prop:cons-bianchi}
Suppose the Einstein equation \eqref{eq:einstein} holds, with $\Lambda$ constant. Then
\begin{equation}
\label{eq:cons-a}
\nabla^{\mu}T^{\mathrm{tot}}_{\mu\nu}=0 .
\end{equation}
\end{proposition}

\begin{proof}
Take the covariant divergence of \eqref{eq:einstein}. The contracted second Bianchi identity gives
$\nabla^{\mu}G_{\mu\nu}=0$ \cite{Wald}, and metric compatibility \eqref{eq:LC} with constancy of
$\Lambda$ gives $\nabla^{\mu}(\Lambda g_{\mu\nu})=0$. Hence
$8\pi G_N\,\nabla^\mu T^{\mathrm{tot}}_{\mu\nu}=0$, and \eqref{eq:cons-a} follows since
$G_N\neq0$.
\end{proof}

This argument uses the Einstein equation: it is valid on a full solution of the coupled system,
including the gravitational field equation, and it does not establish \eqref{eq:cons-a} from the
gauge and matter equations alone. Route (b) below derives the same conclusion without the Einstein
equation, as a Noether identity of diffeomorphism invariance.

\subsection{Route (b): diffeomorphism invariance, independent of the Einstein equation}

\begin{theorem}[Diffeomorphism invariance and on-shell covariant conservation]
\label{thm:cons}
Let
$S_{\mathrm{ng}}:=S_{\mathrm{YM}}+S_{\mathrm{m}}+S_{\mathrm{rel}}$
denote the nongravitational part of the action \eqref{eq:Stot}, and assume:
\begin{enumerate}[label=(C\arabic*),leftmargin=2.8em]
\item\label{C1} \emph{(Invariance.)} $S_{\mathrm{ng}}$ is invariant under every smooth
      one-parameter family of allowed arrows $(\hat F_t,f_t)$ of $\Ggeo$ with compactly supported
      generator $\xi$, the fields $(g,A,\Phi)$ transforming as in \eqref{eq:delta-xi}. This is the
      hypothesis \ref{R2} of Definition~\ref{def:action}, restricted to arrows near the identity;
\item\label{C2} \emph{(Invariance of the fixed data.)} all data held fixed during the variation are
      themselves invariant under that flow. Concretely: any background or external structure in
      $\mathcal{B}$ entering $S_{\mathrm{ng}}$ is either dragged along with the fields or is
      preserved by the flow; and every relation coordinate $r_\alpha$ in regime (b) of
      Convention~\ref{conv:two-regimes} (a fixed topological or boundary sector label) satisfies
      $\delta_\xi r_\alpha=0$, which holds because such labels are $\Ggeo$-invariant by
      construction and $\xi$ generates arrows of $\Ggeo$;
\item\label{C3} \emph{(Boundary terms.)} $\xi$ has compact support in the interior of $M$, or the
      boundary terms generated by the integrations by parts below vanish under the treatment
      declared in Assumption~\ref{ass:variational};
\item\label{C4} \emph{(Full Euler coefficients.)} the coefficients of $\delta_{\hat\xi} A_\mu$ and
      $\delta_{\hat\xi}\Phi$ in $\delta_\xi S_{\mathrm{ng}}$ are the \emph{full} Euler expressions
      appearing on the left sides of \eqref{eq:ym} and \eqref{eq:matter}---that is, they include the
      YM, matter, \emph{and} relation contributions, computed with the chain rule
      \eqref{eq:chain-rule};
\item\label{C5} \emph{(Availability of lifts.)} for \emph{every} smooth vector field $\xi$ with
      compact support in the interior of $M$, the admissibility class and allowed-arrow structure of
      $\Ggeo$ contain a smooth allowed one-parameter family of arrows $(\hat F_t,f_t)$ whose base
      generator is $\xi$, with the fields $(g,A,\Phi)$ and the fixed data of $\mathcal{B}$
      transformed as in \eqref{eq:delta-xi}. Thus every such $\xi$ generates an allowed flow to which
      \ref{C1} and \ref{C2} apply.
\end{enumerate}
Write the full Euler operators of the gauge and matter sectors as
\begin{equation}
\label{eq:euler-ops}
\mathcal{G}^{\nu}
:=\frac{1}{g_{\mathrm{YM}}^{2}}\,D_\mu F^{\mu\nu}
-\bigl(J^{\nu}_{\mathrm{m}}+J^{\nu}_{\mathrm{rel}}\bigr),
\qquad
\mathcal{E}_\Phi
:=\mathcal{E}^{\mathrm{m}}_{\Phi}+\mathcal{E}^{\mathrm{rel}}_{\Phi},
\end{equation}
so that the gauge equation \eqref{eq:ym} reads $\mathcal{G}^{\nu}=0$ and the matter equation
\eqref{eq:matter} reads $\mathcal{E}_\Phi=0$.

Then diffeomorphism invariance gives the integrated variational identity displayed in the proof. On
every configuration satisfying $\mathcal{G}^{\nu}=0$ and $\mathcal{E}_\Phi=0$, but not necessarily
the Einstein equation, one has
\begin{equation}
\label{eq:cons-b}
\nabla^{\mu}T^{\mathrm{tot}}_{\mu\nu}=0 .
\end{equation}
If the sectors exchange energy--momentum, the individual divergences $\nabla^\mu T^{X}_{\mu\nu}$
need not vanish; only the sum does.
\end{theorem}

\begin{proof}
By \ref{C1} and \ref{C2}, $\delta_\xi S_{\mathrm{ng}}=0$ \emph{identically} in the fields---this is
the off-shell content of diffeomorphism invariance, and \ref{C2} is exactly what licenses it, since
a non-invariant fixed datum would contribute an uncancelled term. Expanding the total variation in
the three field arguments, with the full Euler coefficients of \ref{C4},
\begin{equation}
\label{eq:ward-expand}
0=\delta_\xi S_{\mathrm{ng}}
=\int_M\left[
\frac{\delta S_{\mathrm{ng}}}{\delta g^{\mu\nu}}\,\delta_\xi g^{\mu\nu}
+\Bigl\langle \frac{\delta S_{\mathrm{ng}}}{\delta A_\mu},\,\delta_{\hat\xi} A_\mu\Bigr\rangle
+\Bigl\langle\!\Bigl\langle \frac{\delta S_{\mathrm{ng}}}{\delta \Phi},\,\delta_{\hat\xi} \Phi
\Bigr\rangle\!\Bigr\rangle
\right]\mathrm{d}^4x ,
\end{equation}
We first identify the three coefficients in \eqref{eq:ward-expand}. By \eqref{eq:Jdef},
\eqref{eq:Edef} and Lemma~\ref{lem:YMvar}, the coefficient of $\delta_{\hat\xi} A_\mu$ is $\sqrt{-g}$
times $\mathcal{G}^{\mu}$, and the coefficient of $\delta_{\hat\xi}\Phi$ is $\sqrt{-g}$ times
$\mathcal{E}_\Phi$, in the notation \eqref{eq:euler-ops}. The coefficient of $\delta_\xi g^{\mu\nu}$
is, by \eqref{eq:Tdef} and \eqref{eq:Ttot}, equal to $-\tfrac12\sqrt{-g}\,T^{\mathrm{tot}}_{\mu\nu}$.
Hence \eqref{eq:ward-expand} reads
\begin{equation}
\label{eq:ward-three-terms}
0=
-\frac12\int_M T^{\mathrm{tot}}_{\mu\nu}\,\delta_\xi g^{\mu\nu}\,\dvol
+\int_M\Bigl[
\bigl\langle \mathcal{G}^{\mu},\,\delta_{\hat\xi}A_\mu\bigr\rangle
+\bigl\langle\!\bigl\langle \mathcal{E}_\Phi,\,\delta_{\hat\xi}\Phi\bigr\rangle\!\bigr\rangle
\Bigr]\,\dvol .
\end{equation}
Equation \eqref{eq:ward-three-terms} is the integrated variational identity referred to in the
statement. It holds for every compactly supported $\xi$ generating a family of allowed arrows, and
we now impose the gauge and matter equations on it.

Accordingly, let the configuration satisfy $\mathcal{G}^{\mu}=0$ and $\mathcal{E}_\Phi=0$. The
second integral in \eqref{eq:ward-three-terms} then vanishes term by term, whatever the variations
$\delta_{\hat\xi}A_\mu$ and $\delta_{\hat\xi}\Phi$ happen to be. This is also the step announced in
Convention~\ref{conv:lie}: whatever lift $\hat\xi$ of $\xi$ was chosen, and hence whatever vertical
terms $D_\mu\varepsilon$ and $-\rho_*(\varepsilon)\Phi$ appear in \eqref{eq:delta-xi-local}, they
enter only through $\delta_{\hat\xi}A_\mu$ and $\delta_{\hat\xi}\Phi$, which are here contracted with
the vanishing $\mathcal{G}^{\mu}$ and $\mathcal{E}_\Phi$; the conclusion is therefore independent of
the choice of lift. Hence \eqref{eq:ward-three-terms} reduces to
\[
0=-\frac12\int_M T^{\mathrm{tot}}_{\mu\nu}\,\delta_\xi g^{\mu\nu}\,\dvol .
\]
Insert $\delta_\xi g^{\mu\nu}=\nabla^{\mu}\xi^{\nu}+\nabla^{\nu}\xi^{\mu}$ from
\eqref{eq:delta-xi-metric}, whose sign is fixed by \eqref{eq:lie-inverse}, and use the symmetry of
$T^{\mathrm{tot}}_{\mu\nu}$:
\[
0=-\int_M T^{\mathrm{tot}}_{\mu\nu}\,\nabla^{\mu}\xi^{\nu}\,\dvol
=\int_M \bigl(\nabla^{\mu}T^{\mathrm{tot}}_{\mu\nu}\bigr)\,\xi^{\nu}\,\dvol ,
\]
the last step by covariant integration by parts, whose boundary term vanishes by \ref{C3}. By
\ref{C5}, every smooth compactly supported vector field $\xi$ in the interior of $M$ arises as the
base generator of an allowed flow, so the identity above holds for all such $\xi$; since $\xi$ is
thereby arbitrary among compactly supported interior fields, the fundamental lemma of the calculus
of variations gives \eqref{eq:cons-b}.

This argument nowhere uses the Einstein equation: \eqref{eq:cons-b} follows from diffeomorphism
invariance \cite{Noether1918} together with the gauge and matter equations alone.
\end{proof}

\subsection{The gauge Ward identity}

\begin{proposition}[Gauge-current identity]
\label{prop:gauge-current}
Assume $S_{\mathrm{m}}+S_{\mathrm{rel}}$ is invariant under the vertical (gauge) subgroup
$\Gau(P)$; that any relation coordinate in regime (b) is gauge invariant (as it is, by
construction); and that Assumption~\ref{ass:variational} holds. Then on solutions of the matter
equation \eqref{eq:matter},
\begin{equation}
\label{eq:current-cons}
D_\mu\bigl(J^{\mu}_{\mathrm{m}}+J^{\mu}_{\mathrm{rel}}\bigr)=0 .
\end{equation}
\end{proposition}

\begin{proof}
Let $\epsilon\in\Gamma(\mathrm{ad}P)$ be a compactly supported infinitesimal gauge parameter, acting
by
\begin{equation}
\label{eq:inf-gauge}
\delta_\epsilon A_\mu=D_\mu\epsilon,
\qquad
\delta_\epsilon\Phi=-\rho_*(\epsilon)\,\Phi ,
\end{equation}
where $\rho_*\colon\gfrak\to\mathfrak{gl}(W)$ is the induced Lie-algebra representation and
$\delta_\epsilon\Phi$ is valued in $W$ (no inner product on $W$ is used). Gauge invariance gives
$\delta_\epsilon(S_{\mathrm{m}}+S_{\mathrm{rel}})=0$ identically:
\[
0=\int_M\left(
\Bigl\langle\frac{\delta (S_{\mathrm{m}}+S_{\mathrm{rel}})}{\delta A_\mu},\;D_\mu\epsilon\Bigr\rangle
+\Bigl\langle\!\Bigl\langle
\frac{\delta (S_{\mathrm{m}}+S_{\mathrm{rel}})}{\delta \Phi},\;-\rho_*(\epsilon)\Phi
\Bigr\rangle\!\Bigr\rangle
\right)\mathrm{d}^4x ,
\]
the second pairing being the duality pairing of Convention~\ref{conv:matter}(i) between the
dual-valued Euler operator and the $W$-valued variation. On the matter equation \eqref{eq:matter} the
second term vanishes. Using \eqref{eq:Jdef}, the first term is
\[
0=-\int_M\bigl\langle J^{\mu}_{\mathrm{m}}+J^{\mu}_{\mathrm{rel}},\;D_\mu\epsilon\bigr\rangle\,\dvol .
\]
Covariant integration by parts---valid because $\langle\cdot,\cdot\rangle$ is $\Ad$-invariant, hence
$D$-compatible, and because the boundary term vanishes by Assumption~\ref{ass:variational}---gives
\[
0=\int_M\bigl\langle D_\mu\bigl(J^{\mu}_{\mathrm{m}}+J^{\mu}_{\mathrm{rel}}\bigr),\;
\epsilon\bigr\rangle\,\dvol .
\]
Since $\epsilon$ is arbitrary and $\langle\cdot,\cdot\rangle$ is nondegenerate,
\eqref{eq:current-cons} follows.
\end{proof}

\begin{remark}[Consistency of the gauge equation with the current identity]
\label{rem:consistency-check}
Identity \eqref{eq:current-cons} is consistent with, and indeed implied on shell by, the gauge
equation \eqref{eq:ym}. Applying $D_\nu$ to \eqref{eq:ym} requires evaluating
$D_\nu D_\mu F^{\mu\nu}$. Decomposing the second covariant derivative into its symmetric and
antisymmetric parts in $(\mu,\nu)$ and contracting with the antisymmetric $F^{\mu\nu}$, only the
antisymmetric part survives:
\[
D_\nu D_\mu F^{\mu\nu}
=\tfrac12\,[D_\nu,D_\mu]\,F^{\mu\nu}
=\tfrac12\,\bigl[F_{\nu\mu},\,F^{\mu\nu}\bigr] ,
\]
by the Ricci identity for the gauge-covariant derivative on a section of $\mathrm{ad}P$, namely
$[D_\nu,D_\mu]X=[F_{\nu\mu},X]$. (The Levi-Civita curvature terms cancel between the two
spacetime indices of the antisymmetric $F^{\mu\nu}$, by the first Bianchi identity for $R$.) The
remaining bracket vanishes: writing $F_{\nu\mu}=-F_{\mu\nu}$ gives
$[F_{\nu\mu},F^{\mu\nu}]=-[F_{\mu\nu},F^{\mu\nu}]$, and $[X,X]=0$ for the $\gfrak$-valued
quantity $X=F_{\mu\nu}$ contracted with its own index-raised counterpart. Hence
$D_\nu D_\mu F^{\mu\nu}=0$, and \eqref{eq:ym} therefore forces
$D_\mu\bigl(J^{\mu}_{\mathrm{m}}+J^{\mu}_{\mathrm{rel}}\bigr)=0$ after relabelling the free index, in
agreement with
Proposition~\ref{prop:gauge-current}. The two derivations are independent: the proposition uses
only gauge invariance and the matter equation, whereas this computation uses the gauge equation.
\end{remark}

\begin{remark}[Ward identities as a constraint on realizations]
\label{rem:noether-constraint}
Theorem~\ref{thm:cons} and Proposition~\ref{prop:gauge-current} function as selection rules. Every
effective representation of a completed source relation must satisfy them,
because they follow from nothing more than the invariance requirement \ref{R2} that made
$S_{\mathrm{rel}}$ a functional of the physical orbit in the first place. Consequently, a proposed
correction to the Tier-1 field equations that violates the relevant Ward identity is \emph{not} a
geometric realization of an invariant completion. This is a sharp, checkable necessary condition on
any candidate coupling of completed relations to Tier-1 dynamics.
\end{remark}

% ======================================================================
\section{From completion to Tier-1 representation}
\label{sec:tier1}

The completed geometric architecture established by Theorems~\ref{thm:descent},
\ref{thm:restricted-realization} and \ref{thm:field-eqs} is the factorization
\begin{equation}
\label{eq:architecture}
\Sgeo
\;\xrightarrow{\;\;\egeo\;\;}\;
\Egeo
\;\xrightarrow{\;\;u_A\;\;}\;
T ,
\qquad
u_A\circ\egeo=\pgeo .
\end{equation}
The first map retains the readout together with the selected invariant geometric relations; the
second forgets those completion coordinates and returns the original Tier-1 state. Neither map
chooses a coordinate chart, a gauge potential, or a preferred representative of a physical orbit:
$\Sgeo$ is already the orbit space, and $\egeo$ is built from $\Ggeo$-invariant functions. This is
what it means for the completion to be geometric rather than notational.

When a relation-dependent action is specified in the sense of Definition~\ref{def:action}, the
completed state enters Tier-1 dynamics through the derived response fields
\begin{equation}
\label{eq:responses-summary}
T^{\mathrm{rel}}_{\mu\nu},
\qquad
J^{\mu}_{\mathrm{rel}},
\qquad
\mathcal{E}^{\mathrm{rel}}_{\Phi} .
\end{equation}
A distinction must be drawn here, since the two halves of the construction carry invariance in
different senses, and conflating them would misdescribe the responses.

\emph{Invariance.} What is invariant are the \emph{scalars} and the \emph{orbit-defined
information}: the action $S_{\mathrm{rel}}$ is a $\Ggeo$-invariant functional by hypothesis
\ref{R2}; the relation coordinates $r_\alpha$ and the readout $\pgeo$ are functions on the orbit
space $\Sgeo$ by Definition~\ref{def:invariant-relation}; and the completion \eqref{eq:architecture}
is assembled from these.

\emph{Covariance.} The responses \eqref{eq:responses-summary}, by contrast, are \emph{not} invariant
pointwise, and are not scalars. They are the functional derivatives of an invariant action with
respect to the fields, hence tensorial objects on $M$ carrying spacetime and representation indices:
$T^{\mathrm{rel}}_{\mu\nu}$ is a symmetric $2$-tensor, $J^{\mu}_{\mathrm{rel}}$ an
$\mathrm{ad}P$-valued spacetime vector field, and $\mathcal{E}^{\mathrm{rel}}_{\Phi}$ a section of
the dual matter bundle. Under an arrow $(\hat F,f)$ of $\Ggeo$ they transform \emph{covariantly}---by pushforward of
tensor indices along $f$ and by the adjoint or dual representation along $\hat F$---rather than
being fixed. Thus, for instance, $J^{\mu}_{\mathrm{rel}}\mapsto \Ad_{h^{-1}}J^{\mu}_{\mathrm{rel}}$
under a gauge transformation $h$, and $T^{\mathrm{rel}}_{\mu\nu}\mapsto (f^{-1})^{*}
T^{\mathrm{rel}}_{\mu\nu}$ under an allowed diffeomorphism $f$.

The correct statement is therefore the following. What descends to the orbit space is the invariant
action and the relation data; what the variation produces is a family of covariant response tensors,
one per gauge and coordinate representative, transforming into one another under the arrows of
$\Ggeo$. It is precisely this covariance---not pointwise invariance---that makes the field equations
\eqref{eq:einstein}--\eqref{eq:matter} well posed as equations between tensors, and that makes the
Ward identities of Section~\ref{sec:noether} meaningful: $\nabla^{\mu}T^{\mathrm{tot}}_{\mu\nu}=0$
is a covariant equation, and is preserved under the arrows of $\Ggeo$ rather than being invariant
under them in the sense of Definition~\ref{def:invariant-relation}. By
Theorem~\ref{thm:field-eqs} and Section~\ref{sec:noether}, the responses
\eqref{eq:responses-summary} are correctly normalized, Ward-consistent, covariant objects obtained
by functional variation of a $\Ggeo$-invariant action.

This is the whole of what the present paper establishes about dynamics, and the boundary of the
argument should be stated precisely. Paper~4 constructs covariant response fields from
$\Ggeo$-invariant actions. It does not determine whether the resulting Tier-1 evolution closes autonomously on $T$, which
completion coordinates (if any) must be retained for a closed description, what happens when
unresolved variables are eliminated, or whether such elimination produces memory or other
history-dependent effective behaviour. In particular, no inference of Tier-1 autonomy may be drawn
merely from the observation that some response functional happens to descend to a function of $t$ at
the kinematic level considered here: descent of a response at one level does not establish closure of
an evolution, and the two questions are not the same. Autonomous closure, retained variables,
elimination, memory, and projected evolution require the separate analysis of Paper~5, and are not
treated here.

% ======================================================================
\section{Conclusion}
\label{sec:conclusion}

Canonical source completion acquires physical meaning only on a domain where its relation
coordinates are jointly realized by compatible geometric data. Four results establish where that
domain lies and what it supports.

The geometric descent theorem (Theorem~\ref{thm:descent}) constructs the global metric, principal
bundle, connection, and matter section from compatible local data, uniquely up to bundle
isomorphism, with the global time-orientation condition treated explicitly
(Remark~\ref{rem:descent-time}); its corollary (Corollary~\ref{cor:realizability-boundary}) exhibits
the realizability boundary, since relation values admitting no compatible descent data are simply
outside the geometric realization domain. Passage to the orbit space of the bundle-automorphism
groupoid (Definition~\ref{def:sgeo}) removes coordinate and gauge redundancy \emph{before}
completion, so that no gauge representative is taken for a physical state; the isotropy of a
configuration is its stabilizer, which is in general a proper subgroup of the ambient automorphism
group (Remark~\ref{rem:isotropy}). The restricted geometric realization theorem
(Theorem~\ref{thm:restricted-realization}) then identifies the joint image of the readout and the
declared invariant relations with the canonical minimal completion of Paper~3---relative to that
family, and only on that domain---while the orbit reconstruction criterion
(Proposition~\ref{prop:orbit-reconstruction}) reduces exact reconstruction to a separation question
in invariant theory, answered case by case and never by fiat.

The flat-circle classification (Theorem~\ref{thm:hol-class} and
Corollaries~\ref{cor:hol-completion}--\ref{cor:spec}) shows in a fully solvable case that a global
relation can completely parametrize a reduced source even when local curvature vanishes identically,
and that a coarser declared relation leaves a strictly coarser completion behind. Finally, invariant
relation-dependent actions (Theorem~\ref{thm:field-eqs}) convert completed structure into correctly
normalized Einstein, Yang--Mills, and matter equations, with the associated covariant response
fields derived rather than postulated, and the Ward--Noether identities (Proposition~\ref{prop:cons-bianchi},
Theorem~\ref{thm:cons}, Proposition~\ref{prop:gauge-current}) enforce covariant conservation and
supply a sharp necessary condition on any candidate realization.

This supplies the geometric bridge from abstract completion to observable Tier-1 structure. It does
not assert that every abstract completion is geometrically realizable, that any finite family of
local invariants separates Lorentzian geometries, or anything about the physical status of our
universe. The dynamical questions of closure, retention, elimination, and memory belong to Paper~5.

% ======================================================================
\begin{thebibliography}{99}

\bibitem{AharonovBohm}
Y.~Aharonov and D.~Bohm,
\emph{Significance of Electromagnetic Potentials in the Quantum Theory},
Physical Review \textbf{115} (1959), 485--491.

\bibitem{AmbroseSinger}
W.~Ambrose and I.~M.~Singer,
\emph{A theorem on holonomy},
Transactions of the American Mathematical Society \textbf{75} (1953), 428--443.

\bibitem{Bleecker}
D.~Bleecker,
\emph{Gauge Theory and Variational Principles},
Addison--Wesley, Reading, MA, 1981.

\bibitem{ChernWeil}
S.-S.~Chern,
\emph{Complex Manifolds Without Potential Theory},
2nd ed., Springer, New York, 1979.

\bibitem{GibbonsHawking1977}
G.~W.~Gibbons and S.~W.~Hawking,
\emph{Action integrals and partition functions in quantum gravity},
Physical Review D \textbf{15} (1977), 2752--2756.

\bibitem{KobayashiNomizu1}
S.~Kobayashi and K.~Nomizu,
\emph{Foundations of Differential Geometry}, Volume I,
Interscience Publishers, New York, 1963.

\bibitem{KobayashiNomizu2}
S.~Kobayashi and K.~Nomizu,
\emph{Foundations of Differential Geometry}, Volume II,
Interscience Publishers, New York, 1969.

\bibitem{MilnorStasheff}
J.~W.~Milnor and J.~D.~Stasheff,
\emph{Characteristic Classes},
Annals of Mathematics Studies 76, Princeton University Press, Princeton, NJ, 1974.

\bibitem{Noether1918}
E.~Noether,
\emph{Invariante Variationsprobleme},
Nachrichten von der Gesellschaft der Wissenschaften zu G\"ottingen,
Mathematisch-Physikalische Klasse (1918), 235--257.

\bibitem{Pravda2002}
V.~Pravda, A.~Pravdov\'a, A.~Coley, and R.~Milson,
\emph{All spacetimes with vanishing curvature invariants},
Classical and Quantum Gravity \textbf{19} (2002), 6213--6236.

\bibitem{Steenrod}
N.~Steenrod,
\emph{The Topology of Fibre Bundles},
Princeton University Press, Princeton, NJ, 1951.

\bibitem{Wald}
R.~M.~Wald,
\emph{General Relativity},
University of Chicago Press, Chicago, 1984.

\bibitem{WuYang}
T.~T.~Wu and C.~N.~Yang,
\emph{Concept of nonintegrable phase factors and global formulation of gauge fields},
Physical Review D \textbf{12} (1975), 3845--3857.

\bibitem{YangMills1954}
C.~N.~Yang and R.~L.~Mills,
\emph{Conservation of Isotopic Spin and Isotopic Gauge Invariance},
Physical Review \textbf{96} (1954), 191--195.

\bibitem{York1972}
J.~W.~York, Jr.,
\emph{Role of Conformal Three-Geometry in the Dynamics of Gravitation},
Physical Review Letters \textbf{28} (1972), 1082--1085.

\bibitem{ST1}
J.~Rodgers,
\emph{Source--Readout Non-Equivalence: Descent and Equivariant Reconstruction Obstructions},
Zenodo, 2026.
\href{https://doi.org/10.5281/zenodo.21369967}{doi:10.5281/zenodo.21369967}.

\bibitem{ST2}
J.~Rodgers,
\emph{Target-Relative Necessity of Completion: When Readout Loss Obstructs, and What a Sufficient
Extension Must Retain},
Zenodo, 2026.
\href{https://doi.org/10.5281/zenodo.21370501}{doi:10.5281/zenodo.21370501}.

\bibitem{ST3}
J.~Rodgers,
\emph{Canonical Minimal Source Completion: The Coarsest Readout Extension on which a Nominated
Family of Source Relations Becomes Well Defined},
Zenodo, 2026.
\href{https://doi.org/10.5281/zenodo.21370763}{doi:10.5281/zenodo.21370763}.

\end{thebibliography}

\end{document}
